\documentclass[10pt,a4paper]{article}
\usepackage[margin=2.2cm]{geometry}
\usepackage{fontspec}
\setmonofont{DejaVu Sans Mono}[Scale=0.84]
\newfontfamily\unifont{DejaVu Serif}[Scale=0.9]
\usepackage{amsmath,amssymb,mathtools}
\setcounter{MaxMatrixCols}{20}
\usepackage{longtable,booktabs,array}
\usepackage{fancyvrb}
\usepackage[dvipsnames]{xcolor}
\usepackage[colorlinks=true,linkcolor=NavyBlue,urlcolor=NavyBlue]{hyperref}
\usepackage{microtype}
\usepackage{needspace}
\usepackage{tocloft}\setlength{\cftbeforesecskip}{0.45em}
\setlength{\parindent}{0pt}\setlength{\parskip}{0.55em}
\setlength{\emergencystretch}{8em}
\makeatletter
\renewcommand\section{\@startsection{section}{1}{\z@}{-3.5ex \@plus -1ex \@minus -.2ex}{2.3ex \@plus.2ex}{\normalfont\Large\bfseries\raggedright}}
\renewcommand\subsection{\@startsection{subsection}{2}{\z@}{-3.25ex\@plus -1ex \@minus -.2ex}{1.5ex \@plus .2ex}{\normalfont\large\bfseries\raggedright}}
\renewcommand\subsubsection{\@startsection{subsubsection}{3}{\z@}{-3.25ex\@plus -1ex \@minus -.2ex}{1.5ex \@plus .2ex}{\normalfont\normalsize\bfseries\raggedright}}
\makeatother
\renewcommand{\arraystretch}{1.15}
\allowdisplaybreaks
\begin{document}
\title{The radion of 7+1 spacetime\\[0.4em]\large The size of the four extra dimensions as a field in ordinary spacetime — on the 4-torus and on the 4-sphere; a guide for students, with worked examples checked in section 26 of the notebook 7+1.nb and by an independent Python program, and a check of the author’s note on the radion}
\author{Repository: github.com/once-ere/7plus1 — prepared for Nash\\ Notebook: 7+1.nb (Wolfram Language 15.0.1), evaluated Wed 28 Oct 2026 08:57:47\\ Document generated by tools/make\_docs.py on 2026-10-28}
\date{}
\maketitle
\begin{abstract}
The four extra space-like coordinates $x_4,\dots,x_7$ of the 7+1 model describe a small compact space. Its overall SIZE can vary from place to place and from time to time in ordinary spacetime; seen from four dimensions this size is a scalar field, the RADION. This document defines the radion precisely — through the Einstein-frame metric $ds^2=e^{-4\beta(x)}g_{E\mu\nu}dx^\mu dx^\nu-e^{2\beta(x)}\gamma_{ij}dy^idy^j$, whose field $\beta=\ln(b/b_0)$ has the canonical form $\phi=\sqrt{12}\,M_{\rm Pl}\beta$ — and works out, for two choices of the extra space, the 4-TORUS $T^4$ of 7+1.nb and the 4-SPHERE $S^4$, everything that follows from it: the reduced action, the potential and the mass of the radion (the torus is a modulus, the sphere is not), how the radion couples to the world brane and to the Higgs boson, the fifth force and the Eöt-Wash bound, the Kaluza–Klein levels and the Dirac zero modes of dirac16complex on the two spaces, and the contrast with the warped radion of Randall and Sundrum. It checks, statement by statement, a note of the author on the radion and corrects what the checks require (a factor 2 in the canonical normalization, the Einstein-frame factor of a bulk potential, the signs of the brane tensions, the normalization constant and the mass of the warped radion, the five-dimensional equations offered for eight dimensions). Every equation used is written here. Section 26 of 7+1.nb contains 33 checks (33 True); the independent program tools/ref\_radion.py contains 31 (31 True); the 15 groups of numbers printed by both agree to a largest relative difference of $8.2\times10^{-11}$.
\end{abstract}
\tableofcontents
\Needspace{12\baselineskip}\section*{1. Results at a glance}\addcontentsline{toc}{section}{1. Results at a glance}
{\small
\begin{longtable}{>{\raggedright\arraybackslash}p{\dimexpr 0.2500\linewidth-2\tabcolsep\relax}>{\raggedright\arraybackslash}p{\dimexpr 0.6200\linewidth-2\tabcolsep\relax}>{\raggedright\arraybackslash}p{\dimexpr 0.1300\linewidth-2\tabcolsep\relax}}
\toprule
\textbf{question} & \textbf{answer} & \textbf{where} \\
\midrule\endfirsthead
\toprule
\textbf{question} & \textbf{answer} & \textbf{where} \\
\midrule\endhead
\endfoot
\bottomrule\endlastfoot
What is the radion? & the field $\beta(x)=\ln(b/b_0)$ (torus) or $\ln(r/r_0)$ (sphere) that measures the size of the extra space at the point $x$ of ordinary spacetime; canonical form $\phi=\sqrt{12}\,M_{\rm Pl}\beta$ & §4 \\
Why the factor $e^{-4\beta}$ in front of $g_E$? & it keeps the coefficient of the four-dimensional curvature constant (EINSTEIN FRAME): $e^{2\cdot(-2\beta)}e^{4\beta}=1$; in general $2\alpha+n\beta=0$ & §4.2 \\
What is its kinetic term? & $6M_{\rm Pl}^2(\partial\beta)^2=\frac12(\partial\phi)^2$; for $n$ extra dimensions $\frac{M_{\rm Pl}^2}2\,\frac{n(n+2)}2\,(\partial\beta)^2$ — coefficients 1.50, 4.00, 7.50, 12.0 for $n=1,\dots,4$ & §4.3, §4.5 \\
What is the physical volume? & $V=V_0e^{4\beta}=V_0e^{2\phi/(\sqrt3M_{\rm Pl})}$; $V_0=b_0^4L_c^4$ ($T^4$) or $\frac{8\pi^2}3r_0^4$ ($S^4$) & §4.4 \\
Does it have a potential? & on $T^4$ with $\Lambda=0$: no — a MODULUS; on $S^4$ the curvature gives $-6M_{\rm Pl}^2e^{-6\beta}/r_0^2$; the bulk cosmological constant gives $M_{\rm Pl}^2\Lambda e^{-4\beta}$ & §6 \\
How is it stabilized, and how heavy is it? & $T^4$: $\Lambda+\kappa U_0=0$, $U_1=0$, $m_r^2=\kappa K_U/12$. $S^4$: $(\Lambda+\kappa U_0)r_0^2=6$, $\kappa U_1r_0^2=-12$, $m_r^2=\kappa K_U/12-2/r_0^2$ & §6.3–6.4 \\
Can the sphere be held by $\Lambda$ alone? & only at a maximum: $dS_4\times S^4$ with $\Lambda r_0^2=9$, $H^2r_0^2=1$, $m^2=-6\,H^2$ (unstable); a 4-form flux can hold it: $\Lambda r_0^2=4.50$, $\kappa\rho_Fr_0^2=1.50$, $m^2r_0^2=6$ & §6.5–6.6 \\
How does it couple to the world brane? & through $h_{\mu\nu}=e^{-4\beta}g_{E\mu\nu}$, on $T^4$ and $S^4$ alike: $\mathcal L=\frac{\phi}{\sqrt3M_{\rm Pl}}T^\mu{}_\mu$; a massless radion adds a fifth force of strength $\frac23$ & §7.1–7.2 \\
How heavy must it be? & Eöt-Wash: heavier than about $\hbar c/38.6\,\mu$m $=5.112$ meV; the radion of 7+1.nb ($m_r=10$ GeV) has the range $1.973\times10^{-17}$ m & §7.2 \\
Does it mix with the Higgs boson? & yes, through $\xi R(h)H^\dagger H$ on the brane, with $\gamma=v/(\sqrt3M_{\rm Pl})=5.838\times10^{-17}$: for $\xi=1$, $m_{r0}=10$ GeV the mixing angle is $-3.53\times10^{-16}$ and the Higgs rates change by $3.4\times10^{-31}$ — unobservable & §7.3 \\
Torus or sphere — what changes? & levels $2\pi|\vec n|/\ell_\perp$ versus $\sqrt{l(l+3)}/r$; the brane couples to $r_4(k)$ states per torus level and to ONE zonal state per sphere level; dirac16complex has zero modes on $T^4$ and none on $S^4$ (lowest $|\lambda|=2/r$) & §8 \\
Flat or warped? & in 7+1 (unwarped) the radion couples with $1/(\sqrt3M_{\rm Pl})$; in Randall–Sundrum with $1/\langle\mu\rangle$, $1.41\times10^{15}$ times stronger for $\langle\mu\rangle=3$ TeV & §9 \\
What does the check of the note find? & its ideas hold; the checks correct a factor 2 in the canonical normalization, a missing Einstein-frame factor, the sign of a junction condition (and so of the brane tensions), $f_\mu$ and the warped radion mass, mixed normalizations, five-dimensional equations used in eight dimensions, and an incomplete kinetic Lagrangian of the Higgs–radion mixing & §11 \\*
\end{longtable}}
\Needspace{12\baselineskip}\section*{2. Conventions}\addcontentsline{toc}{section}{2. Conventions}
\begin{itemize}
\item Coordinates $X=(x_0,\dots,x_7)$: $x_0$ is time — the direction of the octonion unit $e_0=1$; $x_1,x_2,x_3$ are ordinary space; $x_4,x_5,x_6,x_7$ are the four extra (auxiliary) space-like coordinates. Ordinary spacetime points are written $x=(x_0,\dots,x_3)$, points of the extra space $y=(x_4,\dots,x_7)$; Greek indices $\mu,\nu$ run over $0,\dots,3$, Latin $i,j$ over $4,\dots,7$, capital $A,B$ over $0,\dots,7$.
\item The flat metric is $\eta=\mathrm{diag}(+1,-1,\dots,-1)$: simply the squares of the octonion basis, $e_0^2=+1$, $e_k^2=-1$. The signature is "mostly minus"; xAct declares the spacetime with the signature list $\{7,1,0\}$: seven space-like and one time-like dimension.
\item Curvature: $R^\rho{}_{\sigma\mu\nu}=\partial_\mu\Gamma^\rho{}_{\nu\sigma}-\partial_\nu\Gamma^\rho{}_{\mu\sigma}+\Gamma^\rho{}_{\mu\lambda}\Gamma^\lambda{}_{\nu\sigma}-\Gamma^\rho{}_{\nu\lambda}\Gamma^\lambda{}_{\mu\sigma}$, $R_{\sigma\nu}=R^\rho{}_{\sigma\rho\nu}$, $R=g^{\sigma\nu}R_{\sigma\nu}$, $G^A{}_B=R^A{}_B-\frac12\delta^A_BR$. In this signature a round sphere in the space-like block has NEGATIVE $R$ (for example $R=-12/r_0^2$ for the 4-sphere of radius $r_0$); the sign of $R$ flips with the sign of the metric.
\item Field equations and action: $G^A{}_B-\Lambda\,\delta^A_B=\kappa\,T^A{}_B$, $\kappa=8\pi G_8=M_8^{-6}$, from
\end{itemize}
\begin{equation*}
S=-\frac1{2\kappa}\int d^8X\sqrt{|G|}\,\big(R+2\Lambda\big)-\int d^8X\sqrt{|G|}\,U+S_{\rm brane}+S_{\rm matter},
\end{equation*}
where $\Lambda$ is the bulk cosmological constant ($\Lambda>0$: de Sitter), $U$ the stabilizing energy density of the bulk (section 6), $S_{\rm brane}$ the action of the world brane (section 7). Units $\hbar=c=1$; $M_{\rm Pl}=2.435\times10^{18}$ GeV is the reduced Planck mass, $8\pi G_4=M_{\rm Pl}^{-2}$; the examples use $M_8=10$ TeV and $V_\perp=M_{\rm Pl}^2/M_8^6$ (the volume of the extra space today, section 4.4).

\textbf{The author’s note.} The note checked here is an e-mail of 28 October 2026 in which the author collected a discussion of the radion: the Kaluza–Klein ansatz with the factors $e^{2\alpha\phi}$, $e^{2\beta\phi}$, the Einstein frame and the canonical normalization, the physical volume, the radion as a modulus and the fifth force, stabilization (Goldberger–Wise), the warped alternative of Randall and Sundrum (RS1) with its junction conditions, normalization and mass, the mixing of the radion with the Higgs boson, and its effects on Higgs decays. The note uses the signature $(-,+,+,+)$, in which the scalar curvature has the opposite sign; its statements are translated into the conventions above. Section 11 lists every statement with the verdict of the checks.

\Needspace{12\baselineskip}\section*{3. The idea: a breathing extra space}\addcontentsline{toc}{section}{3. The idea: a breathing extra space}
Imagine at every point of ordinary space a tiny four-dimensional space, too small to see. If it were rigid, physics in ordinary space would never notice its size, only its effects on the constants of nature (for example $G_4=G_8/V$). If it can change size — bigger here, smaller there, oscillating in time — the size becomes a FIELD in ordinary spacetime. A field that changes from point to point carries kinetic energy, $\frac12(\partial\phi)^2$, so it behaves like a particle: the radion. Two questions then decide everything: how is the size defined so that its kinetic energy has the standard form, and what energy does the extra space store as a function of its size (the potential, whose curvature is the mass of the radion)?

Two extra spaces are used as examples throughout:

\begin{enumerate}
\item \textbf{The 4-torus $T^4$} (the space of 7+1.nb): $x_4,\dots,x_7$ periodic with period $L_c$, the flat metric $b^2\delta_{ij}$, size $b$, proper volume $V=b^4L_c^4$. It is flat, so its size costs no curvature energy.
\item \textbf{The 4-sphere $S^4$} (for comparison): $x_4,x_5,x_6\in[0,\pi]$, $x_7\in[0,2\pi)$ angles, radius $r$, the metric $r^2\hat\gamma$ with $\hat\gamma=\mathrm{diag}(1,\sin^2x_4,\sin^2x_4\sin^2x_5,\sin^2x_4\sin^2x_5\sin^2x_6)$, proper volume $V=\frac{8\pi^2}3r^4$. It is curved, so its size costs energy — a sphere cannot be a modulus.
\end{enumerate}
\Needspace{12\baselineskip}\section*{4. The definition}\addcontentsline{toc}{section}{4. The definition}
\Needspace{8\baselineskip}\subsection*{4.1 The metric ansatz}\addcontentsline{toc}{subsection}{4.1 The metric ansatz}
The RADION $\beta(x)$ is defined by the eight-dimensional metric
\begin{equation*}
ds^2=e^{-4\beta(x)}\,g_{E\mu\nu}(x)\,dx^\mu dx^\nu-e^{2\beta(x)}\,\gamma_{ij}(y)\,dy^idy^j,
\end{equation*}
with a FIXED reference metric $\gamma$ of the extra space — $\gamma=b_0^2\delta_{ij}$ on $T^4$, $\gamma=r_0^2\hat\gamma$ on $S^4$ — and the four-dimensional EINSTEIN-FRAME metric $g_E$. The size is $b=b_0e^\beta$ on $T^4$ and the radius $r=r_0e^\beta$ on $S^4$: $\beta=\ln(b/b_0)$, the same field as in section 20.3 of 7+1.nb. In the notation of the note, $ds^2=e^{2\alpha\phi}g+e^{2\beta\phi}\gamma$ (signature $(-,+,+,+)$) with constants $\alpha,\beta$ and a field $\phi$: our $\beta(x)$ is the note’s $\beta\phi(x)$ and our $-2\beta(x)$ its $\alpha\phi(x)$.

\Needspace{8\baselineskip}\subsection*{4.2 Why the factor \texorpdfstring{$e^{-4\beta}$}{e-4}: the Einstein frame}\addcontentsline{toc}{subsection}{4.2 Why the factor \texorpdfstring{$e^{-4\beta}$}{e-4}: the Einstein frame}
Under a constant rescaling $g\to\Omega^2g$ the scalar curvature scales as $R\to\Omega^{-2}R$ and the volume element as $\sqrt{|g|}\to\Omega^4\sqrt{|g|}$ (four dimensions). For the 4D block $e^{2\alpha}g$ and the internal block $e^{2\beta}\gamma$ ($n$ extra dimensions) the eight-dimensional volume element and the 4D part of the curvature give
\begin{equation*}
\sqrt{|G|}\,R[G]\supset e^{4\alpha}e^{n\beta}\sqrt{|g|}\sqrt\gamma\cdot e^{-2\alpha}R[g]=e^{2\alpha+n\beta}\sqrt{|g|}\sqrt\gamma\,R[g] .
\end{equation*}
The coefficient of $\sqrt{|g|}R[g]$ — the strength of four-dimensional gravity — is constant if and only if
\begin{equation*}
2\alpha+n\beta=0:\qquad n=4:\ \alpha=-2\beta .
\end{equation*}
This is the Einstein frame. If the factor were absent ($\alpha=0$), Newton’s constant would be $G_4=G_8/V_0e^{4\beta}$, a field — the description known as the Jordan frame; both describe the same physics, but only in the Einstein frame is the graviton cleanly separated from the radion.

\Needspace{8\baselineskip}\subsection*{4.3 The reduced action and the canonical radion}\addcontentsline{toc}{subsection}{4.3 The reduced action and the canonical radion}
Insert the ansatz into the action of section 2 and integrate over the extra space ($\int d^4y\sqrt\gamma=V_0$). Section 26.1 of 7+1.nb proves, for the homogeneous metrics $g_E=\mathrm{diag}(N^2,-a^2,-a^2,-a^2)$ with arbitrary $N(x_0)$, $a(x_0)$, $\beta(x_0)$ (the Euler–Lagrange expressions of the difference of the two sides vanish identically, so the difference is a total derivative), that
\begin{equation*}
-\frac1{2\kappa}\int d^4y\sqrt{|G|}\,(R+2\Lambda)-\int d^4y\sqrt{|G|}\,U=\sqrt{|g_E|}\,\Big[-\frac{M_{\rm Pl}^2}2R_E+6M_{\rm Pl}^2\,g_E^{\mu\nu}\partial_\mu\beta\,\partial_\nu\beta-V(\beta)\Big]+\text{(total derivative)},
\end{equation*}
\begin{equation*}
M_{\rm Pl}^2=\frac{V_0}\kappa,\qquad V(\beta)=e^{-4\beta}\Big[M_{\rm Pl}^2\Lambda+V_0U-\frac{M_{\rm Pl}^2}2R_{\rm int}\Big],\qquad R_{\rm int}=\begin{cases}0&(T^4)\\[2pt] \dfrac{12}{r_0^2e^{2\beta}}=\dfrac{12}{r^2}&(S^4)\end{cases}
\end{equation*}
($R_{\rm int}$ is the scalar curvature of the extra space itself, positive for the sphere; $R_E$ that of $g_E$ in the signature of section 2). The kinetic term $6M_{\rm Pl}^2(\partial\beta)^2$ is that of a canonically normalized scalar,
\begin{equation*}
6M_{\rm Pl}^2(\partial\beta)^2=\frac12(\partial\phi)^2\iff \phi=\sqrt{12}\,M_{\rm Pl}\,\beta=\sqrt{12}\,M_{\rm Pl}\ln\frac b{b_0},
\end{equation*}
the CANONICAL RADION. Its mass is the curvature of the potential at the minimum, $m_r^2=d^2V/d\phi^2=V''(\beta)/(12M_{\rm Pl}^2)$.

\Needspace{6\baselineskip}\subsubsection*{Worked example 1 (the torus): the reduction step by step}
On $T^4$ the internal metric $e^{2\beta}b_0^2\delta_{ij}$ is flat, so $R_{\rm int}=0$, and $\sqrt{|G|}=e^{-8\beta}e^{4\beta}\sqrt{|g_E|}b_0^4=e^{-4\beta}\sqrt{|g_E|}b_0^4$. The eight-dimensional curvature of the ansatz contains $e^{4\beta}R_E$ (from the 4D block) plus terms with $\partial\beta$; after the integration over $T^4$ ($V_0=b_0^4L_c^4$) and integrations by parts only $6M_{\rm Pl}^2(\partial\beta)^2$ survives, and the potential is $V=e^{-4\beta}(M_{\rm Pl}^2\Lambda+V_0U)$. With $\Lambda=U=0$ the torus is a modulus: any constant $\beta$ is a solution.

\Needspace{6\baselineskip}\subsubsection*{Worked example 2 (the sphere): the extra potential}
On $S^4$ the internal metric $e^{2\beta}r_0^2\hat\gamma$ has $R_{ij}=3\hat\gamma_{ij}$ and $R_{\rm int}=12/(r_0^2e^{2\beta})$ (section 5.2). The term $-\frac1{2\kappa}\sqrt{|G|}R$ contributes $+\frac{V_0}{2\kappa}e^{-4\beta}R_{\rm int}$ to the Lagrangian, i.e. the potential
\begin{equation*}
V_{\rm curv}(\beta)=-\frac{M_{\rm Pl}^2}2e^{-4\beta}\frac{12}{r_0^2e^{2\beta}}=-\frac{6M_{\rm Pl}^2}{r_0^2}\,e^{-6\beta}:
\end{equation*}
negative and growing without bound as the sphere shrinks ($\beta\to-\infty$). Without anything else the sphere collapses; its size must be held (section 6).

\Needspace{8\baselineskip}\subsection*{4.4 The physical volume and the size}\addcontentsline{toc}{subsection}{4.4 The physical volume and the size}
The proper volume of the extra space at the point $x$ is
\begin{equation*}
V(x)=\int d^4y\,\sqrt{\det\big(e^{2\beta(x)}\gamma_{ij}\big)}=V_0\,e^{4\beta(x)}=V_0\exp\Big(\frac{2\,\phi(x)}{\sqrt3\,M_{\rm Pl}}\Big),\qquad V_0=\begin{cases}b_0^4L_c^4&(T^4)\\ \frac{8\pi^2}3r_0^4&(S^4)\end{cases}
\end{equation*}
Measuring $V$ measures the radion. Today, with $M_8=10$ TeV, the volume is fixed by Newton’s constant: $V_\perp=M_{\rm Pl}^2/M_8^6=5.929\times10^{12}$ GeV$^{-4}$, the torus has the proper size $\ell_\perp=V_\perp^{1/4}=1560.4$ GeV$^{-1}$ $=3.079\times10^{-13}$ m, and a sphere of the same volume has the radius $r_{S^4}=(3V_\perp/(8\pi^2))^{1/4}=688.94$ GeV$^{-1}$ $=1.3595\times10^{-13}$ m.

\Needspace{8\baselineskip}\subsection*{4.5 \texorpdfstring{$n$}{n} extra dimensions, and the factor of the note}\addcontentsline{toc}{subsection}{4.5 \texorpdfstring{$n$}{n} extra dimensions, and the factor of the note}
For a torus $T^n$ ($D=4+n$) with the Einstein-frame factor $e^{-n\beta}$ in front of $g_E$, section 26.1 SOLVES the reduction for the coefficient of the kinetic term: with $\beta=\phi$,
\begin{align*}
-\frac1{2\kappa}\int d^ny\sqrt{|G|}\,R&=\sqrt{|g_E|}\,\Big[-\frac{M_{\rm Pl}^2}2R_E+\frac{M_{\rm Pl}^2}2\,c_n\,(\partial\beta)^2\Big]+\text{(total derivative)}, \\
c_n&=\frac{n(n+2)}2:\qquad c_1=1.50,\quad c_2=4.00,\quad c_3=7.50,\quad c_4=12.0 .
\end{align*}
With $\beta=\beta_n\varphi$ the kinetic term is $\frac{M_{\rm Pl}^2}2c_n\beta_n^2(\partial\varphi)^2$. The note chooses
\begin{equation*}
\alpha=-\sqrt{\frac n{2(n+2)}},\qquad \beta_n=\sqrt{\frac2{n(n+2)}}\qquad(2\alpha+n\beta_n=0\ \checkmark),
\end{equation*}
and states the reduced action $\frac{V_n}{2\kappa^2}\int\sqrt{-g}\,[R-\frac12(\partial\varphi)^2]$ (signature $(-,+,+,+)$). The check gives $c_n\beta_n^2=1$ for the note’s $\beta_n$ ($\beta_4^2=0.08333$): the reduced action is $\frac{V_n}{2\kappa^2}\int\sqrt{-g}\,[R-(\partial\varphi)^2]$ — \textbf{a factor 2 in the kinetic term}. The form with $\frac12$ needs $c_n\beta_n^2=\frac12$:
\begin{equation*}
\beta_n=\frac1{\sqrt{n(n+2)}},\qquad \alpha=-\frac n2\beta_n=-\sqrt{\frac n{4(n+2)}}\qquad(n=4:\ \beta_4^2=0.04167).
\end{equation*}
For the eight-dimensional case the note uses $\alpha=-1/\sqrt3$, $\beta=1/\sqrt{12}$ with the action $\int\sqrt{-g}[\frac{M_{\rm Pl}^2}2R-\frac12(\partial\phi)^2]$. This is right if the exponent is read as $\beta\phi/M_{\rm Pl}$: then $\beta(x)_{\rm here}=\phi/(\sqrt{12}M_{\rm Pl})$, exactly the canonical radion of section 4.3, and $V=V_0e^{4\beta\phi/M_{\rm Pl}}=V_0e^{(2/\sqrt3)\phi/M_{\rm Pl}}$ as the note says.

\Needspace{6\baselineskip}\subsubsection*{Worked example 3: one extra dimension (\texorpdfstring{$n=1$}{n=1}), the classic Kaluza–Klein scalar}
For $n=1$: $c_1=\frac32$, so the canonical choice is $\beta_1^2=\frac13$, $\alpha^2=\frac1{12}$ — the textbook values of five-dimensional Kaluza–Klein theory ($ds^2=e^{2\alpha\varphi}g+e^{2\beta\varphi}dz^2$ with $\beta=-2\alpha$, $\alpha^2=\frac1{12}$). The note’s formula would give $\beta_1^2=\frac23$, twice the textbook value.

\Needspace{12\baselineskip}\section*{5. The two extra spaces}\addcontentsline{toc}{section}{5. The two extra spaces}
\Needspace{8\baselineskip}\subsection*{5.1 The 4-torus}\addcontentsline{toc}{subsection}{5.1 The 4-torus}
$T^4=\mathbb R^4/(L_c\mathbb Z)^4$ with the metric $b^2\delta_{ij}$: flat ($R_{ijkl}=0$), volume $b^4L_c^4$, isometries the four translations and the 384 signed permutations of the axes. Only the product $bL_c$ is physical. The torus is the auxiliary space of 7+1.nb.

\Needspace{8\baselineskip}\subsection*{5.2 The 4-sphere}\addcontentsline{toc}{subsection}{5.2 The 4-sphere}
The unit 4-sphere is the set of points of $\mathbb R^5$ at distance 1 from the origin, with the coordinates
\begin{equation*}
(X_1,\dots,X_5)=r\big(\sin x_4\sin x_5\sin x_6\cos x_7,\ \sin x_4\sin x_5\sin x_6\sin x_7,\ \sin x_4\sin x_5\cos x_6,\ \sin x_4\cos x_5,\ \cos x_4\big);
\end{equation*}
the embedding induces the metric $r^2\hat\gamma$, $\hat\gamma=\mathrm{diag}(1,\sin^2x_4,\sin^2x_4\sin^2x_5,\sin^2x_4\sin^2x_5\sin^2x_6)$, and
\begin{equation*}
R_{ij}=3\,\hat\gamma_{ij},\qquad R_{\rm int}=\frac{12}{r^2},\qquad V=r^4\int_0^\pi\!\!\int_0^\pi\!\!\int_0^\pi\!\!\int_0^{2\pi}\sin^3x_4\sin^2x_5\sin x_6\,dx_4dx_5dx_6dx_7=\frac{8\pi^2}3r^4,
\end{equation*}
all checked in section 26.1 (in the signature of 7+1.nb the space-like block $-r^2\hat\gamma$ has $R=-12/r^2$). Its isometries are the 10 rotations of $\mathbb R^5$, the group $SO(5)$; the north pole is $x_4=0$, the point $(0,0,0,0,r)$.

\Needspace{8\baselineskip}\subsection*{5.3 Side by side}\addcontentsline{toc}{subsection}{5.3 Side by side}
{\small
\begin{longtable}{>{\raggedright\arraybackslash}p{\dimexpr 0.3000\linewidth-2\tabcolsep\relax}>{\raggedright\arraybackslash}p{\dimexpr 0.3500\linewidth-2\tabcolsep\relax}>{\raggedright\arraybackslash}p{\dimexpr 0.3500\linewidth-2\tabcolsep\relax}}
\toprule
\textbf{} & \textbf{$T^4$} & \textbf{$S^4$} \\
\midrule\endfirsthead
\toprule
\textbf{} & \textbf{$T^4$} & \textbf{$S^4$} \\
\midrule\endhead
\endfoot
\bottomrule\endlastfoot
reference metric $\gamma$ & $b_0^2\delta_{ij}$ & $r_0^2\hat\gamma$ \\
size & $b=b_0e^\beta$ & $r=r_0e^\beta$ \\
volume $V_0e^{4\beta}$ & $V_0=b_0^4L_c^4$ & $V_0=\frac{8\pi^2}3r_0^4$ \\
internal curvature $R_{\rm int}$ & $0$ & $12/r^2$ \\
curvature potential & none & $-6M_{\rm Pl}^2e^{-6\beta}/r_0^2$ \\
isometries & $U(1)^4$ translations, 384 signed permutations & $SO(5)$ (10 rotations) \\
radius today ($V=V_\perp$) & $\ell_\perp=3.079\times10^{-13}$ m & $r_{S^4}=1.359\times10^{-13}$ m \\
lowest Kaluza–Klein mass (scalar) & $2\pi/\ell_\perp=4.0265$ MeV & $2/r_{S^4}=2.9030$ MeV \\
Dirac zero modes & 4 (constant spinors) & none (lowest $|\lambda|=2/r$) \\*
\end{longtable}}
\Needspace{12\baselineskip}\section*{6. The potential and the mass of the radion}\addcontentsline{toc}{section}{6. The potential and the mass of the radion}
\Needspace{8\baselineskip}\subsection*{6.1 How each energy scales with the size}\addcontentsline{toc}{subsection}{6.1 How each energy scales with the size}
In the Einstein frame every energy of the bulk carries the factor $\sqrt{|G|}=e^{-4\beta}\sqrt{|g_E|}\sqrt\gamma$ and, if it contains derivatives along the extra space or its curvature, factors $e^{-2\beta}$ from $G^{ij}=-e^{-2\beta}\gamma^{ij}$:

{\small
\begin{longtable}{>{\raggedright\arraybackslash}p{\dimexpr 0.4000\linewidth-2\tabcolsep\relax}>{\raggedright\arraybackslash}p{\dimexpr 0.4000\linewidth-2\tabcolsep\relax}>{\raggedright\arraybackslash}p{\dimexpr 0.2000\linewidth-2\tabcolsep\relax}}
\toprule
\textbf{energy in the bulk} & \textbf{Einstein-frame potential} & \textbf{scaling} \\
\midrule\endfirsthead
\toprule
\textbf{energy in the bulk} & \textbf{Einstein-frame potential} & \textbf{scaling} \\
\midrule\endhead
\endfoot
\bottomrule\endlastfoot
bulk cosmological constant $\Lambda$ & $M_{\rm Pl}^2\Lambda e^{-4\beta}$ & $e^{-4\beta}$ \\
stabilizing energy density $U(\beta)$ & $V_0U\,e^{-4\beta}$ & $e^{-4\beta}$ \\
curvature of $S^4$ & $-6M_{\rm Pl}^2e^{-6\beta}/r_0^2$ & $e^{-6\beta}$ \\
gradient energy of a field along $y$ & $\frac12\int\sqrt\gamma\,|\nabla\Phi|_\gamma^2\,e^{-6\beta}$ & $e^{-6\beta}$ \\
mass energy of that field & $\frac12\int\sqrt\gamma\,m^2\Phi^2\,e^{-4\beta}$ & $e^{-4\beta}$ \\
magnetic 4-form flux through $S^4$ (illustration) & $M_{\rm Pl}^2\kappa\rho_F^{(0)}e^{-12\beta}$ & $e^{-12\beta}$ \\*
\end{longtable}}
With the stabilizing energy density of section 20.3, written around the minimum,
\begin{equation*}
U(\beta)=U_0+U_1\beta+\frac{K_U}2\beta^2,\qquad V(\beta)=M_{\rm Pl}^2\Big[e^{-4\beta}\big(\Lambda+\kappa U(\beta)\big)-\frac{6\,\sigma}{r_0^2}e^{-6\beta}+F\,e^{-12\beta}\Big]
\end{equation*}
($\sigma=0$ for $T^4$, $\sigma=1$ for $S^4$; $F=\kappa\rho_F^{(0)}$ the flux term, zero unless stated). A static FLAT four-dimensional vacuum at $\beta=0$ needs $V(0)=0$ (no four-dimensional cosmological constant) and $V'(0)=0$ (a stationary point); it is stable if $V''(0)>0$, and then $m_r^2=V''(0)/(12M_{\rm Pl}^2)$. Each case below is checked in section 26.2 in two independent ways — from $V(\beta)$ and from the eight-dimensional field equations (section 6.7).

\Needspace{8\baselineskip}\subsection*{6.2 Without stabilization: a modulus, and the fifth force}\addcontentsline{toc}{subsection}{6.2 Without stabilization: a modulus, and the fifth force}
On $T^4$ with $\Lambda=U=0$, $V\equiv0$: the radion is a MODULUS, a massless field whose value is not fixed. A massless radion couples to every mass on the brane (section 7.1) and mediates a long-range force of $\frac23$ the strength of gravity — excluded by tests of gravity. With $\Lambda\ne0$ and $U=0$, $V=M_{\rm Pl}^2\Lambda e^{-4\beta}$ has no minimum: the torus runs away. On $S^4$ the curvature term alone drives the sphere to collapse. A realistic model therefore needs a STABILIZER.

\Needspace{8\baselineskip}\subsection*{6.3 Example 1: the torus with \texorpdfstring{$U$}{U}}\addcontentsline{toc}{subsection}{6.3 Example 1: the torus with \texorpdfstring{$U$}{U}}
\begin{equation*}
V(0)=M_{\rm Pl}^2(\Lambda+\kappa U_0)=0,\qquad V'(0)=M_{\rm Pl}^2\big(-4(\Lambda+\kappa U_0)+\kappa U_1\big)=0,\qquad V''(0)=M_{\rm Pl}^2\kappa K_U,
\end{equation*}
\begin{equation*}
\Lambda+\kappa U_0=0,\qquad U_1=0,\qquad m_r^2=\frac{\kappa K_U}{12}\quad\text{(as in section 20.3 of 7+1.nb).}
\end{equation*}
\Needspace{8\baselineskip}\subsection*{6.4 Example 2: the sphere with \texorpdfstring{$U$}{U}}\addcontentsline{toc}{subsection}{6.4 Example 2: the sphere with \texorpdfstring{$U$}{U}}
The curvature must be balanced:
\begin{align*}
V(0)&=M_{\rm Pl}^2\Big(\Lambda+\kappa U_0-\frac6{r_0^2}\Big)=0,\qquad V'(0)=M_{\rm Pl}^2\Big(-4(\Lambda+\kappa U_0)+\kappa U_1+\frac{36}{r_0^2}\Big)=0, \\
V''(0)&=M_{\rm Pl}^2\Big(16(\Lambda+\kappa U_0)-8\kappa U_1+\kappa K_U-\frac{216}{r_0^2}\Big),
\end{align*}
\begin{equation*}
(\Lambda+\kappa U_0)\,r_0^2=6,\qquad \kappa U_1\,r_0^2=-12,\qquad m_r^2=\frac{\kappa K_U}{12}-2\,\frac1{r_0^2}.
\end{equation*}
A sphere needs a stiffer stabilizer than a torus: its curvature lowers the mass by $2/r_0^2$. With $r_0=r_{S^4}$ today, $2/r_0^2$ corresponds to a mass of $\sqrt2/r_{S^4}=2.053$ MeV — negligible against a radion mass of 10 GeV.

\Needspace{8\baselineskip}\subsection*{6.5 Example 3: the sphere held by \texorpdfstring{$\Lambda$}{} alone — unstable}\addcontentsline{toc}{subsection}{6.5 Example 3: the sphere held by \texorpdfstring{$\Lambda$}{} alone — unstable}
With $U=0$, $V=M_{\rm Pl}^2(\Lambda e^{-4\beta}-6e^{-6\beta}/r_0^2)$ has a stationary point at $V'(0)=0$, i.e. $\Lambda r_0^2=9$, where $V(0)=3M_{\rm Pl}^2/r_0^2>0$: four-dimensional de Sitter space with $H^2=V(0)/(3M_{\rm Pl}^2)$, $H^2r_0^2=1$ — the eight-dimensional solution $dS_4\times S^4$, which section 26.2 also verifies in the eight-dimensional field equations. But $V''(0)=-72M_{\rm Pl}^2/r_0^2$:
\begin{equation*}
m^2=\frac{V''(0)}{12M_{\rm Pl}^2}=-\frac6{r_0^2}=-6\,H^2<0 .
\end{equation*}
The radion sits at a MAXIMUM (a tachyon): a sphere held only by a positive cosmological constant either collapses or decompactifies.

\Needspace{8\baselineskip}\subsection*{6.6 Example 4: the sphere held by a flux (illustration)}\addcontentsline{toc}{subsection}{6.6 Example 4: the sphere held by a flux (illustration)}
7+1.nb has no 4-form field, but the classic way to hold a sphere is a magnetic 4-form flux through it (Freund–Rubin): its energy density is $\rho_F=q^2/(2b^8)$ for the sphere of radius $b$, with the pressures $p_a=-\rho_F$ along ordinary space and $p_b=+\rho_F$ along the sphere (from its comoving energy $W=a^3b^4q^2/(2b^8)$, as in section 20 of 7+1.nb), and in the Einstein frame it scales as $e^{-12\beta}$. A Minkowski vacuum then exists and is stable:
\begin{equation*}
V(0)=0,\ V'(0)=0:\qquad \Lambda r_0^2=4.50,\qquad \kappa\rho_Fr_0^2=1.50,\qquad m^2r_0^2=6 .
\end{equation*}
\Needspace{8\baselineskip}\subsection*{6.7 The same results from the eight-dimensional field equations}\addcontentsline{toc}{subsection}{6.7 The same results from the eight-dimensional field equations}
For the cosmological metric $ds^2=dx_0^2-a^2d\vec x^2-b^2d\Omega_4^2$ (the sphere of radius $b(x_0)$) the Einstein tensor differs from that of the torus ($ds^2=dx_0^2-a^2d\vec x^2-b^2d\vec y^2$) by
\begin{equation*}
G^A{}_A\big|_{S^4}-G^A{}_A\big|_{T^4}=\frac{(6,6,6,6,3,3,3,3)}{b^2},\qquad R\big|_{S^4}-R\big|_{T^4}=-\frac{12}{b^2},
\end{equation*}
and the radion equation of section 20.3 (Einstein gravity, from the $11$ and $44$ field equations with the $00$ constraint; $H_a=a'/a$, $H_b=b'/b=\beta'$; $\varepsilon$, $p_a$, $p_b$ the energy density and pressures of the bulk matter) becomes
\begin{equation*}
\beta''+3H_a\beta'+4\beta'^2=\kappa\Big(\frac\varepsilon6-\frac{p_a}2+\frac{p_b}3\Big)+\frac\Lambda3-\frac{3\sigma}{b^2}\qquad(\sigma=0:\ T^4,\ \sigma=1:\ S^4).
\end{equation*}
The stabilizer contributes $\varepsilon=U$, $p_a=-U$, $p_b=-U-\frac b4U'(b)$; the static solutions and the linearized restoring terms of this equation reproduce every result of sections 6.3–6.6 (section 26.2).

\Needspace{8\baselineskip}\subsection*{6.8 A bulk field as the stabilizer: the Einstein-frame factor}\addcontentsline{toc}{subsection}{6.8 A bulk field as the stabilizer: the Einstein-frame factor}
The note stabilizes the radion with a bulk scalar $\Phi(y)$ of mass $m_8$ (Goldberger and Wise) and writes its energy, with the note’s $\beta$, as $V(\phi)=\int d^4y\sqrt\gamma\,e^{4\beta\phi}[\frac12e^{-2\beta\phi}(\nabla\Phi)^2+\frac12m_8^2\Phi^2]$. The factor $e^{4\beta\phi}$ is the internal volume element only; the four-dimensional block contributes $e^{4\alpha\phi}=e^{-8\beta\phi}$. Section 26.2 checks, for a static profile $\Phi(x_4,x_5)$ on $S^4$,
\begin{equation*}
\sqrt{|G|}\Big(\frac12G^{AB}\partial_A\Phi\partial_B\Phi-\frac12m_8^2\Phi^2\Big)=-\sqrt{|g_E|}\sqrt\gamma\,e^{-4\beta}\Big(\frac12e^{-2\beta}|\nabla\Phi|^2_\gamma+\frac12m_8^2\Phi^2\Big),
\end{equation*}
so the Einstein-frame potential is $e^{-4\beta}\int\sqrt\gamma\,[\frac12e^{-2\beta}|\nabla\Phi|^2+\frac12m_8^2\Phi^2]$: the gradient energy scales as $e^{-6\beta}$, like the curvature of the sphere, the mass energy as $e^{-4\beta}$, like $U$ — not as $e^{+2\beta}$ and $e^{+4\beta}$. Its minimum, where the two compete, fixes the size. (On $T^4$ a non-constant $\Phi$ needs sources, for example boundary values on branes; the note’s boundary values $v_A$, $v_B$ play this role in the warped case of section 9.)

\Needspace{12\baselineskip}\section*{7. The radion and the world brane}\addcontentsline{toc}{section}{7. The radion and the world brane}
\Needspace{8\baselineskip}\subsection*{7.1 The coupling: the same on \texorpdfstring{$T^4$}{T4} and \texorpdfstring{$S^4$}{S4}}\addcontentsline{toc}{subsection}{7.1 The coupling: the same on \texorpdfstring{$T^4$}{T4} and \texorpdfstring{$S^4$}{S4}}
The world brane is the 3+1-dimensional surface $y=0$ (on $T^4$) or the north pole $x_4=0$ (on $S^4$). Brane fields feel the INDUCED metric, the 4D block of the ansatz at that point:
\begin{equation*}
h_{\mu\nu}(x)=e^{-4\beta(x)}\,g_{E\mu\nu}(x)\qquad\text{(on both spaces: the factor depends only on }n=4\text{).}
\end{equation*}
A particle of mass $M$ at rest on the brane has the action $-M\int d\tau=-M\int dx_0\,e^{-2\beta}=-M\int dx_0\,(1-2\beta+\dots)$, so with $\beta=\phi/(\sqrt{12}M_{\rm Pl})$:
\begin{equation*}
\mathcal L_{\rm int}=\frac{\phi}{\sqrt3\,M_{\rm Pl}}\,T^\mu{}_\mu\qquad(T^\mu{}_\mu=M\ \text{for the particle}),
\end{equation*}
the coupling of section 20.3 of 7+1.nb. For $n$ extra dimensions $h=e^{-n\beta}g_E$ and the canonical $\beta=\phi/(M_{\rm Pl}\sqrt{n(n+2)/2})$ give the coupling $\sqrt{n/(2(n+2))}/M_{\rm Pl}$ (section 26.3).

\Needspace{8\baselineskip}\subsection*{7.2 The fifth force and the Eöt-Wash bound}\addcontentsline{toc}{subsection}{7.2 The fifth force and the Eöt-Wash bound}
A scalar coupled with $g\,\phi M$ to two masses adds the potential $-g^2MM'e^{-m_rr}/(4\pi r)=-\alpha\,G_4MM'e^{-m_rr}/r$ with $\alpha=2(gM_{\rm Pl})^2$. For the radion $g=1/(\sqrt3M_{\rm Pl})$:
\begin{equation*}
\alpha=2\Big(\frac{M_{\rm Pl}}{\sqrt3M_{\rm Pl}}\Big)^2=\frac23,\qquad\text{in general}\quad \alpha=\frac n{n+2} .
\end{equation*}
Without a mass the long-range force would be $(1+\frac23)G_4MM'/r=\frac53G_4MM'/r$ — the factor $\frac53$ that the sum over the images of a mass on the torus produces (section 25.4 of 7+1.nb). The Eöt-Wash torsion balance (Lee et al. 2020) excludes Yukawa forces of gravitational strength ($|\alpha|=1$) with a range $\lambda>38.6\,\mu$m; a radion of strength $\frac23$ must therefore be heavier than about
\begin{equation*}
m_r\gtrsim\frac{\hbar c}{38.6\,\mu{\rm m}}=0.005112\ {\rm eV}.
\end{equation*}
The radion of 7+1.nb (sections 22–24) has $m_r=10$ GeV and the range $\hbar c/(10\,{\rm GeV})=1.973\times10^{-17}$ m: far below any test of gravity.

\Needspace{8\baselineskip}\subsection*{7.3 The radion and the Higgs boson: mixing}\addcontentsline{toc}{subsection}{7.3 The radion and the Higgs boson: mixing}
The Higgs doublet $H$ lives on the brane ($H^\dagger H=\varphi_H^2/2$ in unitary gauge, $\varphi_H=v+h$, $v=246$ GeV). Besides its kinetic term and potential it may couple to the curvature of the induced metric (Giudice, Rattazzi and Wells); written in the signature of 7+1.nb, where the note’s $-\xi\int\sqrt{-g}R\,H^\dagger H$ changes the sign of $R$,
\begin{equation*}
\mathcal L=\sqrt{|h|}\Big[\frac12h^{\mu\nu}\partial_\mu\varphi_H\partial_\nu\varphi_H-\frac\lambda4\big(\varphi_H^2-v^2\big)^2+\frac\xi2R(h)\,\varphi_H^2\Big]+\frac12(\partial r)^2-\frac12m_{r0}^2r^2 \qquad(\xi=\tfrac16:\ \text{conformal}).
\end{equation*}
Here $h=\Omega^2\eta$ with $\Omega=1+r/\Lambda_r+O(r^2)$, $r$ the canonically normalized radion and $\Lambda_r$ its coupling scale: in 7+1 $\Omega=e^{-2\beta}=e^{-\phi/(\sqrt3M_{\rm Pl})}$, $\Lambda_r=-\sqrt3M_{\rm Pl}$ (the sign can be absorbed into $r\to-r$); in Randall–Sundrum $\Lambda_r=\langle\mu\rangle$ (section 9). Define $\gamma=v/\Lambda_r$. Section 26.6 proves:

\begin{enumerate}
\item For $h=\Omega^2\eta$ in four dimensions, $\sqrt{|h|}R(h)=c\,\Omega\,\Box\Omega$ with $c=-6$, $\Box=\partial_0^2-\nabla^2$.
\item Expanded to second order in $(h,r)$ and up to total derivatives (the terms linear in the fields are total derivatives — the vacuum is a solution):
\end{enumerate}
\begin{equation*}
\mathcal L^{(2)}=\frac12(\partial h)^2+\frac12\big(1+6\xi\gamma^2\big)(\partial r)^2+6\xi\gamma\,\partial h\cdot\partial r-\frac12m_{h0}^2h^2-\frac12m_{r0}^2r^2,\qquad m_{h0}^2=2\lambda v^2 .
\end{equation*}
(The coefficients $K_{hr}/(\xi\gamma)=6$ and $(K_{rr}-1)/(\xi\gamma^2)=6$ are read off the Euler–Lagrange equations of the expanded Lagrangian.) The note writes the kinetic terms without the $6\xi\gamma^2$ of the radion, but its normalization factor contains it:
\begin{align*}
h&=h'-\frac{6\xi\gamma}Z\,r',\qquad r=\frac{r'}Z:\qquad \text{canonical}\iff Z^2=1+6\xi\gamma^2(1-6\xi) \\
&\qquad(\text{with the note's literal kinetic terms: }Z^2=1-36\xi^2\gamma^2).
\end{align*}
The mass matrix of $(h',r')$ is $M'=\begin{pmatrix}m_{h0}^2&-6\xi\gamma m_{h0}^2/Z\\-6\xi\gamma m_{h0}^2/Z&(m_{r0}^2+36\xi^2\gamma^2m_{h0}^2)/Z^2\end{pmatrix}$. With $h'=\cos\theta\,h_m-\sin\theta\,r_m$, $r'=\sin\theta\,h_m+\cos\theta\,r_m$ its off-diagonal element becomes $-\frac12(M'_{11}-M'_{22})\sin2\theta+M'_{12}\cos2\theta$, which vanishes for
\begin{align*}
\tan2\theta&=\frac{2M'_{12}}{M'_{11}-M'_{22}}=\frac{12\gamma\xi Z\,m_{h0}^2}{m_{r0}^2-m_{h0}^2\big(Z^2-36\xi^2\gamma^2\big)}, \\
h&=C_h\,h_m+\dots,\quad r=C_r\,h_m+\dots,\qquad C_h=\cos\theta-\frac{6\xi\gamma}Z\sin\theta,\quad C_r=\frac{\sin\theta}Z,
\end{align*}
the formulas of the note. (The note displays the rotation matrix with the primed and unprimed fields interchanged; with that reading the off-diagonal element is $+\frac12(M'_{11}-M'_{22})\sin2\theta+M'_{12}\cos2\theta$ and $\theta$ would change sign.) The physical Higgs $h_m$ couples to Standard Model fields through $h$ (factor $C_h$) and, through its radion part (factor $C_r$), to the trace of their energy–momentum tensor, including the trace anomaly; the note writes, with its values $b_{\rm EM}=-\frac{11}3$, $F_{\rm SM}=-6.5$, $b_{\rm QCD}=7$, $F_{\rm heavy}=-\frac43$,
\begin{equation*}
\frac{\Gamma(h\to\gamma\gamma)}{\Gamma_{\rm SM}}=\Big|C_h+C_r\gamma\frac{b_{\rm EM}}{F_{\rm SM}}\Big|^2,\qquad \frac{\Gamma(h\to gg)}{\Gamma_{\rm SM}}=\Big|C_h+C_r\gamma\frac{b_{\rm QCD}}{F_{\rm heavy}}\Big|^2 .
\end{equation*}
\Needspace{6\baselineskip}\subsubsection*{Worked example 4: the Higgs boson and the radion in 7+1 (on \texorpdfstring{$T^4$}{T4} or \texorpdfstring{$S^4$}{S4})}
In 7+1, $\gamma=v/(\sqrt3M_{\rm Pl})=5.838\times10^{-17}$. For $\xi=1$, $m_{r0}=10$ GeV, $m_{h0}=125.20$ GeV (60-digit arithmetic in both programs):
\begin{align*}
Z-1&=-5.112\times10^{-32},\qquad \theta=-3.525\times10^{-16}, \\
\frac{\Gamma(h\to\gamma\gamma)}{\Gamma_{\rm SM}}-1&=9.947\times10^{-32},\qquad \frac{\Gamma(h\to gg)}{\Gamma_{\rm SM}}-1=3.388\times10^{-31} .
\end{align*}
The flat radion of 7+1 cannot be seen through the Higgs boson. The bounds the note derives from LHC data ($\langle\mu\rangle>2.5$–$3$ TeV, $-0.1<\xi<0.4$) concern the warped radion, for which $\gamma=v/(3\,{\rm TeV})=0.0821$ (the note’s 0.08) and the same formulas give large effects. They are quoted from the note; this document does not repeat the fit to the LHC data.

\Needspace{12\baselineskip}\section*{8. Kaluza–Klein levels, the brane and dirac16complex on \texorpdfstring{$T^4$}{T4} and \texorpdfstring{$S^4$}{S4}}\addcontentsline{toc}{section}{8. Kaluza–Klein levels, the brane and dirac16complex on \texorpdfstring{$T^4$}{T4} and \texorpdfstring{$S^4$}{S4}}
\Needspace{8\baselineskip}\subsection*{8.1 The levels of a scalar field}\addcontentsline{toc}{subsection}{8.1 The levels of a scalar field}
A bulk field $\Psi(x,y)=\sum_n\psi_n(x)Y_n(y)$ with $-\Delta_\gamma Y_n=\mu_n^2Y_n$ appears in ordinary spacetime as a tower of fields of mass $\mu_n$:

{\small
\begin{longtable}{>{\raggedright\arraybackslash}p{\dimexpr 0.2000\linewidth-2\tabcolsep\relax}>{\raggedright\arraybackslash}p{\dimexpr 0.4000\linewidth-2\tabcolsep\relax}>{\raggedright\arraybackslash}p{\dimexpr 0.4000\linewidth-2\tabcolsep\relax}}
\toprule
\textbf{} & \textbf{$T^4$ (period $\ell_\perp$)} & \textbf{$S^4$ (radius $r$)} \\
\midrule\endfirsthead
\toprule
\textbf{} & \textbf{$T^4$ (period $\ell_\perp$)} & \textbf{$S^4$ (radius $r$)} \\
\midrule\endhead
\endfoot
\bottomrule\endlastfoot
eigenfunctions & $e^{2\pi i\vec n\cdot\vec y/\ell_\perp}/\sqrt V$, $\vec n\in\mathbb Z^4$ & spherical harmonics of degree $l=0,1,2,\dots$ \\
masses & $\mu=2\pi|\vec n|/\ell_\perp$ & $\mu=\sqrt{l(l+3)}/r$ \\
multiplicity & $r_4(k)=8\sigma(k)-32\sigma(k/4)$ for $|\vec n|^2=k$ & $d_l=\frac{(2l+3)(l+1)(l+2)}6=\binom{l+4}4-\binom{l+2}4$ \\
first multiplicities & $r_4(1..4)=8,24,32,24$ & $d_{0..8}=1,5,14,30,55,91,140,204,285$ \\
Weyl’s law & $N(\mu)\approx V\mu^4/(32\pi^2)$ & $N(\mu)\approx V\mu^4/(32\pi^2)=(\mu r)^4/12$ \\*
\end{longtable}}
On $S^4$ the number of states with $l(l+3)\le x^2$ is exact: with $L$ the largest such $l$ and $y=L+2$, $N=\sum_{l\le L}d_l=y^2(y^2-1)/12$. For integer $x$ ($y=x$) the relative deviation from Weyl’s law is exactly $-1/x^2$: $-0.0100$, $-0.000100$, $-1.00\times10^{-6}$ at $x=\mu r=10,100,1000$. Between integers the count jumps by $d_L\approx L^3/3$ at each level, so the deviation oscillates with an amplitude of about $4/x$.

\Needspace{6\baselineskip}\subsubsection*{Worked example 5: the two spaces today}
With the same volume $V_\perp$ (the same Newton constant): torus $\ell_\perp=3.079\times10^{-13}$ m, lowest mass $2\pi/\ell_\perp=4.0265$ MeV; sphere $r_{S^4}=1.3595\times10^{-13}$ m, lowest mass $\sqrt{1\cdot4}/r=2.9030$ MeV. Below 5 GeV the torus has 11733511217161 states (an exact lattice count), the sphere 11737500461050 (exact), Weyl’s law predicts $1.17335\times10^{13}$ for both — the two spaces cannot be told apart by counting heavy levels, only by the lowest ones.

\Needspace{8\baselineskip}\subsection*{8.2 What the brane sees}\addcontentsline{toc}{subsection}{8.2 What the brane sees}
\textbf{On $S^4$} the brane at the north pole sees, at each level $l$, only the ZONAL harmonic (the one that depends on $x_4$ alone); all others vanish there. With the Gegenbauer polynomial $C_l^{(3/2)}$:
\begin{equation*}
Z_l(x_4)=\sqrt{\frac{d_l}{V}}\;\frac{C_l^{(3/2)}(\cos x_4)}{C_l^{(3/2)}(1)},\qquad C_l^{(3/2)}(1)=\frac{(l+1)(l+2)}2,\qquad \frac1{\sin^3x_4}\frac d{dx_4}\Big(\sin^3x_4\frac{dZ_l}{dx_4}\Big)=-l(l+3)Z_l,
\end{equation*}
normalized ($\int_{S^4}Z_l^2=1$) and of value $\sqrt{d_l/V}$ at the brane (section 26.4 checks $l=0,\dots,6$ exactly). So the brane couples to ONE state per level with the strength $\sqrt{d_l/V}$. \textbf{On $T^4$} every mode $e^{2\pi i\vec n\cdot\vec y/\ell_\perp}/\sqrt V$ has the value $1/\sqrt V$ at $y=0$; the brane couples to the combination $\sum_{|\vec n|^2=k}\psi_{\vec n}$, again one state per level, with the strength $\sqrt{r_4(k)/V}$. \textbf{Symmetry:} the sphere has the 10 Killing vectors of $SO(5)$, $L_{ab}=X_a\partial_b-X_b\partial_a$; at the north pole $(0,0,0,0,r)$ the 6 with $a,b\le4$ vanish (the brane keeps $SO(4)$) and 4 move it (broken). On the torus the brane keeps the 384 signed permutations and breaks the 4 translations — which is why Kaluza–Klein momentum is not conserved at the brane in either case.

\Needspace{8\baselineskip}\subsection*{8.3 dirac16complex: zero modes on \texorpdfstring{$T^4$}{T4}, none on \texorpdfstring{$S^4$}{S4}}\addcontentsline{toc}{subsection}{8.3 dirac16complex: zero modes on \texorpdfstring{$T^4$}{T4}, none on \texorpdfstring{$S^4$}{S4}}
The 16 components of dirac16complex are a four-dimensional Dirac spinor times a 4-component spinor of the extra space. Its Kaluza–Klein masses are $\sqrt{m^2+|\lambda|^2}$ with $\lambda$ the eigenvalues of the Dirac operator $\not D=\tilde\gamma^a\nabla_a$ of the extra space ($\tilde\gamma_a=i\rho_{a+3}$ Hermitian, $\tilde\gamma_a\tilde\gamma_b+\tilde\gamma_b\tilde\gamma_a=2\delta_{ab}$, $\rho$ the 4×4 matrices of the brane spinors of 7+1.nb). \textbf{On $T^4$} with periodic spinors the constant spinors are zero modes: for $m=0$, four massless brane Dirac fermions (section 25 of 7+1.nb). \textbf{On $S^4$} there are none: the Killing spinors
\begin{equation*}
\varepsilon_s=e^{s\frac i2x_4\tilde\gamma_1}\,e^{\frac12x_5\tilde\gamma_1\tilde\gamma_2}\,e^{\frac12x_6\tilde\gamma_2\tilde\gamma_3}\,e^{\frac12x_7\tilde\gamma_3\tilde\gamma_4}\,\varepsilon_0,\qquad \nabla_i\varepsilon_s=s\,\frac i{2r}\,e_i{}^a\tilde\gamma_a\varepsilon_s,\qquad \not D\varepsilon_s=s\,\frac{2i}r\,\varepsilon_s\qquad(s=\pm1)
\end{equation*}
(frame $e^a=(r,r\sin x_4,r\sin x_4\sin x_5,r\sin x_4\sin x_5\sin x_6)\,dx^{a+3}$; checked exactly in section 26.4) reach the lowest value allowed by Friedrich’s bound $|\lambda|^2\ge\frac{nR}{4(n-1)}=\frac4{r^2}$. The spectrum is $\lambda=\pm i(l+2)/r$, $l=0,1,2,\dots$, with the degeneracy $4\binom{l+3}3=4,16,40,80,140,224,\dots$ for each sign (Camporesi and Higuchi); the spinor Weyl law $N(|\lambda|\le L/r)\approx L^4/3$ agrees to 0.0020 at $L=1000$. The lightest dirac16complex states on $S^4$ are eight Dirac fermions of mass $\sqrt{m^2+4/r^2}$ — with $r=r_{S^4}$, at least 2.9030 MeV.

\Needspace{12\baselineskip}\section*{9. Flat versus warped: the Randall–Sundrum radion}\addcontentsline{toc}{section}{9. Flat versus warped: the Randall–Sundrum radion}
7+1.nb is UNWARPED: its bulk is a product of ordinary spacetime and the extra space. The note contrasts this with the warped model of Randall and Sundrum (RS1), in FIVE dimensions: $ds^2=e^{-2A(y)}\eta_{\mu\nu}dx^\mu dx^\nu-dy^2$, $y\in[-\pi R,\pi R]$ with $y\simeq-y$ (an orbifold), branes at $y=0$ (UV) and $y=\pi R$ (IR, where the Standard Model lives). Section 26.5 checks it with the conventions of section 2.

\Needspace{8\baselineskip}\subsection*{9.1 The background, the tensions and the junction conditions}\addcontentsline{toc}{subsection}{9.1 The background, the tensions and the junction conditions}
\begin{equation*}
G^\mu{}_\nu=\big(3A''-6A'^2\big)\delta^\mu_\nu,\qquad G^y{}_y=-6A'^2;\qquad A=k|y|:\quad \Lambda=-6k^2\ \text{(anti-de Sitter)}.
\end{equation*}
At the branes the slope $A'$ jumps by $+2k$ (UV: from $-k$ to $+k$) and by $-2k$ (IR), so $A''=\pm2k\,\delta$; a brane of tension $f$ ($T^\mu{}_\nu=f\delta(y)\delta^\mu_\nu$: energy $f$, pressure $-f$) balances the $\delta$ part of $G^\mu{}_\nu$ (coefficient 3) if
\begin{equation*}
3\cdot(\pm2k)=\kappa f:\qquad f_{\rm UV}=+\frac{6k}\kappa,\qquad f_{\rm IR}=-\frac{6k}\kappa\qquad(f\kappa/k=6,\ -6).
\end{equation*}
The UV brane has POSITIVE tension, the result of Randall and Sundrum, $V_{\rm hid}=-V_{\rm vis}=24M^3k$ (with $\kappa=1/(4M^3)$, section 9.2). The same follows from the extrinsic curvature: with $K^\mu{}_\nu=\frac12g^{\mu\alpha}\partial_yg_{\alpha\nu}=-A'\delta^\mu_\nu$ (normal $\partial_y$), the $\delta$-function part of $G^\mu{}_\nu$ at a brane is
\begin{equation*}
[K^\mu{}_\nu]^+_--\delta^\mu_\nu[K]^+_-=-[A']+4[A']=3[A']=6k\quad\text{(UV)},\qquad \text{so}\qquad G^\mu{}_\nu\big|_\delta=\kappa T^\mu{}_\nu\big|_\delta\ \Rightarrow\ f_{\rm UV}=\frac{6k}\kappa .
\end{equation*}
In the signature $(-,+,+,+,+)$ of the note $G^A{}_B$ changes sign, so there the junction condition reads $[K_{\mu\nu}-\hat g_{\mu\nu}K]^+_-=-\kappa^2S_{\mu\nu}$ (section 26.5). The note states the positive UV tension first ($\lambda_A=-\lambda_B=6k/\kappa_8^2$), but its Israel-junction computation uses $[K_{\mu\nu}-\hat g_{\mu\nu}K]^+_-=+\kappa^2S_{\mu\nu}$ with $S_{\mu\nu}=-\lambda\hat g_{\mu\nu}$ and ends with $\lambda_{\rm UV}=-6k/\kappa^2$, $\lambda_{\rm IR}=+6k/\kappa^2$: the wrong sign of the junction condition reverses both tensions. (Its intermediate steps, $K^\pm_{\mu\nu}=\mp k\hat g_{\mu\nu}$ and $[K_{\mu\nu}-\hat g_{\mu\nu}K]^+_-=6k\hat g_{\mu\nu}$, are right.)

\Needspace{8\baselineskip}\subsection*{9.2 The Planck mass, the radion and the normalizations}\addcontentsline{toc}{subsection}{9.2 The Planck mass, the radion and the normalizations}
\begin{equation*}
M_{\rm Pl}^2=\frac1\kappa\int_{-\pi R}^{\pi R}e^{-2k|y|}dy=\frac{1-e^{-2\pi kR}}{\kappa k} .
\end{equation*}
Promoting $R$ to a field, $ds^2=e^{-2kT(x)|y|}g_{\mu\nu}dx^\mu dx^\nu-T(x)^2dy^2$ ($y\in[-\pi,\pi]$), the bulk action gives the kinetic term (section 26.5; a total derivative removed)
\begin{equation*}
\mathcal L_{\rm kin}=\frac3{\kappa k}\,\big(\partial e^{-k\pi T}\big)^2=\frac12(\partial\mu)^2\quad\text{for}\quad \mu=f\,e^{-k\pi T},\qquad f^2=\frac6{\kappa k}\quad(f^2\kappa k=6).
\end{equation*}
Here $1/(2\kappa)$ is the coefficient of $\sqrt{|G|}R$ in the bulk action, $M_{\rm Pl}^2/2$ that of $\sqrt{|g|}R_4$, $\frac12$ that of $(\partial\mu)^2$. The literature and the note use three normalizations:

{\small
\begin{longtable}{>{\raggedright\arraybackslash}p{\dimexpr 0.2200\linewidth-2\tabcolsep\relax}>{\raggedright\arraybackslash}p{\dimexpr 0.1200\linewidth-2\tabcolsep\relax}>{\raggedright\arraybackslash}p{\dimexpr 0.1600\linewidth-2\tabcolsep\relax}>{\raggedright\arraybackslash}p{\dimexpr 0.2800\linewidth-2\tabcolsep\relax}>{\raggedright\arraybackslash}p{\dimexpr 0.2200\linewidth-2\tabcolsep\relax}}
\toprule
\textbf{bulk action} & \textbf{$\kappa$} & \textbf{$f$} & \textbf{$M_{\rm Pl}^2\ (\times(1-e^{-2\pi kR}))$} & \textbf{used by} \\
\midrule\endfirsthead
\toprule
\textbf{bulk action} & \textbf{$\kappa$} & \textbf{$f$} & \textbf{$M_{\rm Pl}^2\ (\times(1-e^{-2\pi kR}))$} & \textbf{used by} \\
\midrule\endhead
\endfoot
\bottomrule\endlastfoot
$-\frac1{2\kappa}\int\sqrt{|G|}R$ & $\kappa$ & $\sqrt{6/(\kappa k)}$ & $1/(\kappa k)$ & this document, 7+1.nb \\
$2M^3\int\sqrt{|G|}R$ (signature $(-,+,\dots)$) & $1/(4M^3)$ & $\sqrt{24M^3/k}$ & $4M^3/k$ (with $\frac{M_{\rm Pl}^2}2R_4$) & Randall–Sundrum, Goldberger–Wise \\
$M_5^3\int\sqrt{|G|}R_5$ (signature $(-,+,\dots)$) & $1/(2M_5^3)$ & $\sqrt{12M_5^3/k}$ & $2M_5^3/k$ & the note’s five-dimensional part \\*
\end{longtable}}
The note’s five-dimensional part combines $M_5^3\int\sqrt{|G|}R_5$ with $f=\sqrt{24M_5^3/k}$ and $M_{\rm Pl}^2=(M_5^3/k)(1-e^{-2k\pi\langle T\rangle})$, formulas of the second row (Randall and Sundrum define $M_{\rm Pl}$ through $2M_{\rm Pl}^2\int\sqrt{|g|}R_4$); in its own normalization they read $f=\sqrt{12M_5^3/k}$ and $M_{\rm Pl}^2=(2M_5^3/k)(1-e^{-2k\pi\langle T\rangle})$. Its later version with $\kappa_8^2$ has the right $M_{\rm Pl}^2$ (and the right kinetic term $\frac1{2\kappa}\frac6k(\partial e^{-k\pi T})^2$, which requires the full circle $y\in[-\pi,\pi]$ although it writes $\int_0^\pi$) but $f=\sqrt{12/(\kappa k)}$ instead of $\sqrt{6/(\kappa k)}$.

\Needspace{8\baselineskip}\subsection*{9.3 Goldberger–Wise stabilization and the mass of the warped radion}\addcontentsline{toc}{subsection}{9.3 Goldberger–Wise stabilization and the mass of the warped radion}
A bulk scalar $\Phi(y)$ of mass $m$ with the brane values $\Phi(0)=v_A$, $\Phi(\pi T)=v_B$ obeys, between the branes ($\sqrt{|G|}=e^{-4ky}$, $G^{yy}=-1$),
\begin{equation*}
\Phi''-4k\Phi'-m^2\Phi=0,\qquad \Phi=e^{2ky}\big(C_1e^{\nu ky}+C_2e^{-\nu ky}\big),\qquad \nu=\sqrt{4+\frac{m^2}{k^2}}=2+\epsilon+O(m^4),\quad \epsilon=\frac{m^2}{4k^2},
\end{equation*}
and Goldberger and Wise find the potential $V\approx k\epsilon v_A^2+4ke^{-4k\pi T}(v_B-v_Ae^{-\epsilon k\pi T})^2$, i.e. with $\mu=fe^{-k\pi T}$
\begin{equation*}
V(\mu)=4k\Big(\frac\mu f\Big)^4\Big[v_B-v_A\Big(\frac\mu f\Big)^\epsilon\Big]^2+{\rm const}.
\end{equation*}
Section 26.5 checks: $V'(\langle\mu\rangle)=0$ at $(\langle\mu\rangle/f)^\epsilon=v_B/v_A$, i.e. $k\pi\langle T\rangle=(4k^2/m^2)\ln(v_A/v_B)$ (the note’s value); $V''(\langle\mu\rangle)=8k\epsilon^2v_B^2\langle\mu\rangle^2/f^4$ (the note’s algebra is right); and with $f^2=6/(\kappa k)$
\begin{equation*}
m_r^2=V''(\langle\mu\rangle)=\frac43\,\kappa k^2\epsilon^2v_B^2\Big(\frac{\langle\mu\rangle}f\Big)^2=\frac{k^2v_B^2\epsilon^2}{3M^3}\,e^{-2k\pi\langle T\rangle}\quad(\kappa=1/(4M^3)),
\end{equation*}
the result of Goldberger and Wise (coefficient 1.333; the note’s $\frac23$, coefficient 0.6667, follows from its $f^2=12/(\kappa k)$). Goldberger and Wise add that a more careful treatment of the $\epsilon$ terms turns $\epsilon^2$ into $\epsilon^{3/2}$.

\Needspace{8\baselineskip}\subsection*{9.4 Couplings: TeV-strong versus Planck-weak}\addcontentsline{toc}{subsection}{9.4 Couplings: TeV-strong versus Planck-weak}
On the IR brane $h=(\mu/f)^2g$: the warped radion couples to brane matter with $1/\langle\mu\rangle$ ($\mathcal L=(r/\langle\mu\rangle)T^\mu{}_\mu$ for $\mu=\langle\mu\rangle+r$), the flat radion of 7+1 with $1/(\sqrt3M_{\rm Pl})$. For $\langle\mu\rangle=3$ TeV the warped coupling is stronger by
\begin{equation*}
\frac{\sqrt3\,M_{\rm Pl}}{3\ {\rm TeV}}=1.4058\times10^{15} .
\end{equation*}
\Needspace{8\baselineskip}\subsection*{9.5 Warping in eight dimensions}\addcontentsline{toc}{subsection}{9.5 Warping in eight dimensions}
The note also offers the equations $3A''+6A'^2=-\kappa_8^2\Lambda_8$, $6A'^2=-\kappa_8^2\Lambda_8$ for an eight-dimensional warped metric with four extra dimensions. These are the five-dimensional equations. For $ds^2=e^{-2A(y)}\eta-dy^2-d\vec z^{\,2}$ (warping along one of the four extra directions, $\vec z$ the other three) section 26.5 finds
\begin{equation*}
G^\mu{}_\nu=3A''-6A'^2,\qquad G^y{}_y=-6A'^2,\qquad G^z{}_z=4A''-10A'^2,
\end{equation*}
and with only a cosmological constant all eight must be equal, which forces $A''=0$ and $A'=0$: there is no warped solution of this form. Warped compactifications with several extra dimensions need other sources (fluxes, several branes); 7+1.nb does not use them.

\Needspace{12\baselineskip}\section*{10. The radion through the three regimes of 7+1}\addcontentsline{toc}{section}{10. The radion through the three regimes of 7+1}
7+1.nb follows the size of the auxiliary torus — the radion — through the history of its universe (sections 14 and 20–23):

\begin{enumerate}
\item \textbf{Primordial (Kasner) regime:} near $x_0=0$ ordinary space expands while the torus contracts; in Einstein gravity $a\propto\tau^{p_1}$, $b\propto\tau^{q_1}$ with $p_1=0.546918$, $q_1=-0.160189$: the radion decreases logarithmically, $\beta=q_1\ln\tau+$const.
\item \textbf{Without a stabilizer} a gas of dirac16complex makes the torus re-expand, $b\propto\tau^{2/5}$ (section 21.3: $\tau H_b=0.399995$): the radion runs to $+\infty$ — decompactification.
\item \textbf{Capture:} with the stabilizer $U$ the radion is caught when the expansion rate falls below $m_r$ and oscillates, $\beta\simeq A\sin(m_r\tau+\vartheta)/(m_r(\tau-\tau_E))$ (section 22.1, fit $A=0.4693$, $\vartheta=2.136$). Its oscillation energy is four-dimensional matter; the radion decays mostly into Kaluza–Klein pairs of dirac16complex (section 22.2).
\item \textbf{Today:} frozen at a displaced minimum, $\beta_{\rm eq}=1.7\times10^{-86}$ from the cosmological constant and matter, $5.6\times10^{-24}$ from a brane tension $f_W=M_8^4$ (section 22.4, $m_r=10$ GeV): $G_4=G_8/V$ is constant to that accuracy.
\item \textbf{Far future:} in the de Sitter phase the radion relaxes to $3m_r^2\beta_{\rm eq}=3H^2-12\alpha_2H^4$, of order $H_\infty^2/m_r^2$ (section 23.2); without the stabilizer eight-dimensional de Sitter space would let the torus inflate.
\end{enumerate}
On a sphere instead of the torus the same history would need the stiffer stabilizer of section 6.4, and the primordial contraction would meet the curvature potential $-6M_{\rm Pl}^2e^{-6\beta}/r_0^2$, which accelerates the collapse.

\Needspace{12\baselineskip}\section*{11. The note, statement by statement}\addcontentsline{toc}{section}{11. The note, statement by statement}
{\footnotesize
\begin{longtable}{>{\raggedright\arraybackslash}p{\dimexpr 0.0550\linewidth-2\tabcolsep\relax}>{\raggedright\arraybackslash}p{\dimexpr 0.4500\linewidth-2\tabcolsep\relax}>{\raggedright\arraybackslash}p{\dimexpr 0.3750\linewidth-2\tabcolsep\relax}>{\raggedright\arraybackslash}p{\dimexpr 0.1200\linewidth-2\tabcolsep\relax}}
\toprule
\textbf{\#} & \textbf{statement of the note (paraphrased)} & \textbf{verdict} & \textbf{here} \\
\midrule\endfirsthead
\toprule
\textbf{\#} & \textbf{statement of the note (paraphrased)} & \textbf{verdict} & \textbf{here} \\
\midrule\endhead
\endfoot
\bottomrule\endlastfoot
1 & Kaluza–Klein ansatz $e^{2\alpha\phi}g+e^{2\beta\phi}\gamma$; Einstein frame fixes $\alpha,\beta$ & correct ($2\alpha+n\beta=0$) & §4.1–4.2 \\
2 & $\alpha=-\sqrt{n/(2(n+2))}$, $\beta=\sqrt{2/(n(n+2))}$ give $\frac{V_n}{2\kappa^2}[R-\frac12(\partial\phi)^2]$ & factor 2: they give $R-(\partial\phi)^2$; for $\frac12$: $\beta=1/\sqrt{n(n+2)}$ & §4.5 \\
3 & $V_{\rm phys}=V_ne^{n\beta\phi}$ & correct & §4.4 \\
4 & without a potential the radion is a modulus; a massless radion mediates a fifth force & correct on $T^4$; on $S^4$ the curvature gives a potential; strength $\frac23$ & §6.2, §7.2 \\
5 & eight dimensions: $\alpha=-1/\sqrt3$, $\beta=1/\sqrt{12}$, $\frac{M_{\rm Pl}^2}2R-\frac12(\partial\phi)^2$, $V=V_0e^{(2/\sqrt3)\phi}$ & correct with $\phi$ in units of $M_{\rm Pl}$ in the exponent & §4.5 \\
6 & Goldberger–Wise in eight dimensions: $V=\int\sqrt\gamma e^{4\beta\phi}[\frac12e^{-2\beta\phi}(\nabla\Phi)^2+\frac12m_8^2\Phi^2]$ & missing factor $e^{-8\beta\phi}$: gradient $e^{-6\beta}$, mass $e^{-4\beta}$ & §6.8 \\
7 & $m_r\sim(m_8/M_8)^2M_{4D}$; force $\propto e^{-m_rr}/r^2$ & estimate not derived; in 7+1 $m_r^2=\kappa K_U/12$ ($T^4$), $\kappa K_U/12-2/r_0^2$ ($S^4$); Yukawa suppression right (the force is $\propto e^{-m_rr}(1+m_rr)/r^2$) & §6.3–6.4, §7.2 \\
8 & RS1: orbifold, warp $e^{-2kT|y|}$, $\langle T\rangle=R$, TeV scale, coupling $1/\langle\mu\rangle$ & correct & §9 \\
9 & RS1: $S_5=M_5^3\int R_5$ with $f=\sqrt{24M_5^3/k}$, $M_{\rm Pl}^2=(M_5^3/k)(1-e^{-2k\pi\langle T\rangle})$ & mixed normalizations: $\sqrt{12M_5^3/k}$, $2M_5^3/k$ in its own & §9.2 \\
10 & warped eight dimensions: $3A''+6A'^2=-\kappa^2\Lambda$, $6A'^2=-\kappa^2\Lambda$, $A=k|y|$, AdS$_8$ & these are five-dimensional; in 8D $G^z{}_z=4A''-10A'^2$, no solution & §9.5 \\
11 & tensions $\lambda_A=-\lambda_B=6k/\kappa_8^2$ & correct (UV positive) & §9.1 \\
12 & Israel junction: $\lambda_{\rm UV}=-6k/\kappa_8^2$, $\lambda_{\rm IR}=+6k/\kappa_8^2$ & signs reversed (the junction condition needs $-\kappa^2S_{\mu\nu}$) & §9.1 \\
13 & $S_{\rm eff}=\frac1{2\kappa}\int[\frac1k(1-e^{-2k\pi T})R_4+\frac6k(\partial e^{-k\pi T})^2]$, $M_{\rm Pl}^2=\frac{1-e^{-2k\pi\langle T\rangle}}{\kappa k}$ & correct over $y\in[-\pi,\pi]$ (the note writes $\int_0^\pi$) & §9.2 \\
14 & $f_\mu=\sqrt{12/(\kappa_8^2k)}$ & $\sqrt{6/(\kappa k)}$ & §9.2 \\
15 & Goldberger–Wise: bulk equation, profile, $\nu\approx2+\epsilon$, $V(T)$, $k\pi\langle T\rangle=(4k^2/m^2)\ln(v_A/v_B)$, $V''$ & correct ($V(T)$ quoted from Goldberger–Wise) & §9.3 \\
16 & $m_r^2=\frac23\kappa k^2\epsilon^2v_B^2(\langle\mu\rangle/f)^2$ & $\frac43$ (from $f$); Goldberger–Wise’s result & §9.3 \\
17 & Higgs–radion: $-\xi\int\sqrt{-g}R\,H^\dagger H$; kinetic mixing $-6\xi\gamma\,\partial h\cdot\partial r$ & correct; the radion kinetic term also gets $6\xi\gamma^2$ & §7.3 \\
18 & $Z^2=1+6\xi\gamma^2(1-6\xi)$, the shift, $\tan2\theta$, $C_h$, $C_r$ & correct; the displayed rotation is the inverse & §7.3 \\
19 & trace-anomaly couplings, $\Gamma(h\to\gamma\gamma)$, $\Gamma(h\to gg)$ ratios & structure correct; values quoted & §7.3 \\
20 & LHC bounds $\langle\mu\rangle>2.5$–$3$ TeV, $-0.1<\xi<0.4$ & quoted, not checked; irrelevant for 7+1 ($\gamma=5.8\times10^{-17}$) & §7.3 \\*
\end{longtable}}
\Needspace{12\baselineskip}\section*{12. Validity and limitations}\addcontentsline{toc}{section}{12. Validity and limitations}
\begin{itemize}
\item The radion here is the overall size only. The shape moduli of $T^4$ (nine), the Kaluza–Klein vectors and the deformations of $S^4$ are not treated; on $S^4$ the $l=2$ deformations are massive by the curvature, on $T^4$ they are moduli.
\item The reduction is checked for homogeneous fields $N(x_0)$, $a(x_0)$, $\beta(x_0)$; covariance extends it to $\beta(x)$ and $g_E(x)$, as the canonical kinetic term $\frac12g_E^{\mu\nu}\partial_\mu\phi\partial_\nu\phi$ requires.
\item The 4-form flux of section 6.6 is an illustration: 7+1.nb has no such field.
\item The stabilizer is the phenomenological $U(b)$ of section 20; the Goldberger–Wise mechanism is analyzed only in the warped five-dimensional model (section 9.3), where it is quoted and its algebra checked.
\item The Higgs–radion analysis is at the quadratic order; the trace-anomaly coefficients and the LHC bounds are the note’s.
\item The numbers today assume $M_8=10$ TeV and $M_{\rm Pl}=2.435\times10^{18}$ GeV.
\end{itemize}
\Needspace{12\baselineskip}\section*{13. Verification}\addcontentsline{toc}{section}{13. Verification}
Every worked example of this document is a check of section 26 of 7+1.nb (Wolfram Language 15.0.1; evaluated in a fresh kernel, 0 message cells) and of the independent Python program tools/ref\_radion.py (sympy, numpy, mpmath; its own curvature routine, Euler–Lagrange operators, lattice sieve and 60-digit arithmetic). The checks of section 26:

{\footnotesize
\begin{longtable}{>{\raggedright\arraybackslash}p{\dimexpr 0.0500\linewidth-2\tabcolsep\relax}>{\raggedright\arraybackslash}p{\dimexpr 0.7700\linewidth-2\tabcolsep\relax}>{\raggedright\arraybackslash}p{\dimexpr 0.0800\linewidth-2\tabcolsep\relax}>{\raggedright\arraybackslash}p{\dimexpr 0.1000\linewidth-2\tabcolsep\relax}}
\toprule
\textbf{\#} & \textbf{check} & \textbf{result} & \textbf{section} \\
\midrule\endfirsthead
\toprule
\textbf{\#} & \textbf{check} & \textbf{result} & \textbf{section} \\
\midrule\endhead
\endfoot
\bottomrule\endlastfoot
1 & \texttt{section 26 engine for diagonal metrics: rdScal and rdEin reproduce curvature[] of section 11 for the cosmological metric (R and every G\textasciicircum{}A\_\allowbreak{}A)} & True & 26.1 \\
2 & \texttt{the 4-sphere of radius r0: Ricci\_\allowbreak{}ij = 3 gamma\_\allowbreak{}ij (gamma the unit-sphere metric), Ricci\_\allowbreak{}ij = 0 for i != j, so R(S\textasciicircum{}4) = 12/\allowbreak{}r0\textasciicircum{}2 (in the signature of 7+1.nb the space-like block gives R = -12/\allowbreak{}r0\textasciicircum{}2, also with curvature[] of section 11); the volume of the unit S\textasciicircum{}4 is 8 Pi\textasciicircum{}2/\allowbreak{}3} & True & 26.1 \\
3 & \texttt{radion on T\textasciicircum{}4: for ds\textasciicircum{}2 = e\textasciicircum{}(-4 beta) (N\textasciicircum{}2 dx0\textasciicircum{}2 - a\textasciicircum{}2 dx.dx) - e\textasciicircum{}(2 beta) b0\textasciicircum{}2 dy.dy the reduced action -(1/\allowbreak{}(2 kappa)) Int(R + 2 Lambda) - Int U equals -(M\_\allowbreak{}Pl\textasciicircum{}2/\allowbreak{}2) R\_\allowbreak{}E + 6 M\_\allowbreak{}Pl\textasciicircum{}2 (d beta)\textasciicircum{}2 - e\textasciicircum{}(-4 beta)(M\_\allowbreak{}Pl\textasciicircum{}2 Lambda + V0 U) up to a total derivative (Euler-Lagrange expressions for N, a, beta vanish), M\_\allowbreak{}Pl\textasciicircum{}2 = V0/\allowbreak{}kappa, V0 = b0\textasciicircum{}4 Lc\textasciicircum{}4} & True & 26.1 \\
4 & \texttt{radion on S\textasciicircum{}4: the same reduction for ds\textasciicircum{}2 = e\textasciicircum{}(-4 beta) g\_\allowbreak{}E - e\textasciicircum{}(2 beta) r0\textasciicircum{}2 gamma (unit-sphere gamma) holds with the additional potential -(M\_\allowbreak{}Pl\textasciicircum{}2/\allowbreak{}2) e\textasciicircum{}(-4 beta) 12/\allowbreak{}(r0\textasciicircum{}2 e\textasciicircum{}(2 beta)) = -6 M\_\allowbreak{}Pl\textasciicircum{}2 e\textasciicircum{}(-6 beta)/\allowbreak{}r0\textasciicircum{}2 from the curvature of the sphere, V0 = (8 Pi\textasciicircum{}2/\allowbreak{}3) r0\textasciicircum{}4; R is independent of the angles} & True & 26.1 \\
5 & \texttt{n extra dimensions (torus T\textasciicircum{}n, D = 4 + n, n = 1..4, Einstein frame e\textasciicircum{}(-n B) g\_\allowbreak{}E): the kinetic term is (M\_\allowbreak{}Pl\textasciicircum{}2/\allowbreak{}2)(n(n+2)/\allowbreak{}2)(dB)\textasciicircum{}2; with the note's B = beta\_\allowbreak{}n phi, beta\_\allowbreak{}n = Sqrt[2/\allowbreak{}(n(n+2))], it is (M\_\allowbreak{}Pl\textasciicircum{}2/\allowbreak{}2)(d phi)\textasciicircum{}2, not (M\_\allowbreak{}Pl\textasciicircum{}2/\allowbreak{}2)(1/\allowbreak{}2)(d phi)\textasciicircum{}2; the form with 1/\allowbreak{}2 needs beta\_\allowbreak{}n = 1/\allowbreak{}Sqrt[n(n+2)]} & True & 26.1 \\
6 & \texttt{the note's eight-dimensional values alpha = -1/\allowbreak{}Sqrt[3], beta = 1/\allowbreak{}Sqrt[12]: 2 alpha + 4 beta == 0 (Einstein frame), 4 beta == 2/\allowbreak{}Sqrt[3] (V = V0 e\textasciicircum{}(4 beta phi)), and Sqrt[12] beta == 1, so M\_\allowbreak{}Pl phi\_\allowbreak{}note = Sqrt[12] M\_\allowbreak{}Pl ln(b/\allowbreak{}b0) is the canonical radion of section 20.3} & True & 26.1 \\
7 & \texttt{the cosmological metric with the sphere of radius b(x0) (ds\textasciicircum{}2 = dx0\textasciicircum{}2 - a\textasciicircum{}2 dx.dx - b\textasciicircum{}2 dOmega\_\allowbreak{}4\textasciicircum{}2): G\textasciicircum{}A\_\allowbreak{}A = (torus value) + (6, 6, 6, 6, 3, 3, 3, 3)/\allowbreak{}b\textasciicircum{}2 and R = (torus value) - 12/\allowbreak{}b\textasciicircum{}2} & True & 26.2 \\
8 & \texttt{radion equation on S\textasciicircum{}4 (Einstein gravity; 11 and 44 equations with the 00 constraint): Hb' + 3 Ha Hb + 4 Hb\textasciicircum{}2 == kappa (eps/\allowbreak{}6 - pa/\allowbreak{}2 + pb/\allowbreak{}3) + Lambda/\allowbreak{}3 - 3/\allowbreak{}b\textasciicircum{}2} & True & 26.2 \\
9 & \texttt{radion on T\textasciicircum{}4 with U = U0 + U1 beta + (K\_\allowbreak{}U/\allowbreak{}2) beta\textasciicircum{}2: a flat vacuum at b = b0 needs Lambda + kappa U0 = 0 and U1 = 0, and m\_\allowbreak{}r\textasciicircum{}2 = V''(beta)/\allowbreak{}(12 M\_\allowbreak{}Pl\textasciicircum{}2) = kappa K\_\allowbreak{}U/\allowbreak{}12 (section 20.3)} & True & 26.2 \\
10 & \texttt{radion on S\textasciicircum{}4 with U: Lambda + kappa U0 == 6/\allowbreak{}r0\textasciicircum{}2, kappa U1 == -12/\allowbreak{}r0\textasciicircum{}2 and m\_\allowbreak{}r\textasciicircum{}2 == kappa K\_\allowbreak{}U/\allowbreak{}12 - 2/\allowbreak{}r0\textasciicircum{}2 from V(beta); the eight-dimensional radion equation gives the same static solution and the same mass} & True & 26.2 \\
11 & \texttt{S\textasciicircum{}4 with Lambda only: dS4 x S\textasciicircum{}4 with r0\textasciicircum{}2 = 9/\allowbreak{}Lambda, H\textasciicircum{}2 = 1/\allowbreak{}r0\textasciicircum{}2 = Lambda/\allowbreak{}9 solves the 00, 11 and 44 equations; it is a MAXIMUM of V: m\textasciicircum{}2 = -6 H\textasciicircum{}2 from V(beta), and m\textasciicircum{}2 = -6/\allowbreak{}r0\textasciicircum{}2 = -6 H\textasciicircum{}2 from the eight-dimensional radion equation} & True & 26.2 \\
12 & \texttt{S\textasciicircum{}4 with a 4-form flux (illustration; W = a\textasciicircum{}3 b\textasciicircum{}4 q\textasciicircum{}2/\allowbreak{}(2 b\textasciicircum{}8) gives eps = rho\_\allowbreak{}F, pa = -rho\_\allowbreak{}F, pb = +rho\_\allowbreak{}F): Minkowski x S\textasciicircum{}4 needs Lambda = 9/\allowbreak{}(2 r0\textasciicircum{}2) and kappa rho\_\allowbreak{}F = 3/\allowbreak{}(2 r0\textasciicircum{}2); m\textasciicircum{}2 = 6/\allowbreak{}r0\textasciicircum{}2 from V(beta) and from the eight-dimensional radion equation} & True & 26.2 \\
13 & \texttt{the note's Goldberger-Wise potential in eight dimensions: for a static profile Phi(y) (here Phi(x4, x5)) Sqrt|G| (1/\allowbreak{}2 G\textasciicircum{}AB dPhi dPhi - 1/\allowbreak{}2 m\textasciicircum{}2 Phi\textasciicircum{}2) == -N a\textasciicircum{}3 e\textasciicircum{}(-4 beta) Sqrt[gamma] (1/\allowbreak{}2 e\textasciicircum{}(-2 beta) |grad Phi|\textasciicircum{}2\_\allowbreak{}gamma + 1/\allowbreak{}2 m\textasciicircum{}2 Phi\textasciicircum{}2): the Einstein-frame factor is e\textasciicircum{}(-4 beta) (e\textasciicircum{}(-4 beta phi) in the note's notation), not e\textasciicircum{}(+4 beta phi)} & True & 26.2 \\
14 & \texttt{the radion couples to a mass M on the brane as (phi/\allowbreak{}(Sqrt[3] M\_\allowbreak{}Pl)) M: -M e\textasciicircum{}(-n beta/\allowbreak{}2) = -M + (n/\allowbreak{}2) M beta + ..., with beta = phi/\allowbreak{}(M\_\allowbreak{}Pl Sqrt[n (n + 2)/\allowbreak{}2]) the coupling is Sqrt[n/\allowbreak{}(2 (n + 2))]/\allowbreak{}M\_\allowbreak{}Pl = 1/\allowbreak{}(Sqrt[3] M\_\allowbreak{}Pl) for n = 4, and the fifth-force strength 2 (coupling M\_\allowbreak{}Pl)\textasciicircum{}2 = n/\allowbreak{}(n + 2) = 2/\allowbreak{}3} & True & 26.3 \\
15 & \texttt{Eot-Wash (|alpha| = 1 excluded for lambda > 38.6 micrometres): a radion of strength 2/\allowbreak{}3 must be heavier than about hbar c/\allowbreak{}38.6 micrometres = 5.1 meV; the 10 GeV radion has the range 2.0 10\textasciicircum{}-17 m; Higgs-radion mixing in 7+1 is governed by gamma = v/\allowbreak{}(Sqrt[3] M\_\allowbreak{}Pl) = 5.8 10\textasciicircum{}-17} & True & 26.3 \\
16 & \texttt{zonal harmonics on S\textasciicircum{}4 (l = 0..6, exact): (1/\allowbreak{}sin\textasciicircum{}3 x4) d/\allowbreak{}dx4 (sin\textasciicircum{}3 x4 d/\allowbreak{}dx4 C\_\allowbreak{}l\textasciicircum{}(3/\allowbreak{}2)(cos x4)) == -l (l + 3) C\_\allowbreak{}l\textasciicircum{}(3/\allowbreak{}2), C\_\allowbreak{}l\textasciicircum{}(3/\allowbreak{}2)(1) == (l + 1)(l + 2)/\allowbreak{}2, and Int over S\textasciicircum{}4 of (C\_\allowbreak{}l/\allowbreak{}C\_\allowbreak{}l(1))\textasciicircum{}2 == V/\allowbreak{}d\_\allowbreak{}l (unit sphere), so the normalized zonal harmonic is Sqrt[d\_\allowbreak{}l/\allowbreak{}V] at the north pole; d\_\allowbreak{}l == Binomial[l + 4, 4] - Binomial[l + 2, 4] for l = 0..20} & True & 26.4 \\
17 & \texttt{Weyl's law on S\textasciicircum{}4: the number of levels with l(l + 3) <= (mu r)\textasciicircum{}2 approaches (mu r)\textasciicircum{}4/\allowbreak{}12 = V mu\textasciicircum{}4/\allowbreak{}(32 Pi\textasciicircum{}2) (V = 8 Pi\textasciicircum{}2 r\textasciicircum{}4/\allowbreak{}3): relative deviation -1.0 10\textasciicircum{}-2, -1.0 10\textasciicircum{}-4, -1.0 10\textasciicircum{}-6 at mu r = 10, 100, 1000 (like -1/\allowbreak{}(mu r)\textasciicircum{}2)} & True & 26.4 \\
18 & \texttt{the brane at the north pole of S\textasciicircum{}4: the embedding (r0 sin x4 sin x5 sin x6 cos x7, ..., r0 cos x4) of S\textasciicircum{}4 in R\textasciicircum{}5 induces r0\textasciicircum{}2 gamma; the north pole x4 = 0 is (0, 0, 0, 0, r0); of the 10 rotations L\_\allowbreak{}ab of R\textasciicircum{}5 (the Killing vectors of S\textasciicircum{}4, SO(5)) the 6 with a, b <= 4 vanish there (SO(4), kept by the brane) and the 4 L\_\allowbreak{}a5 move it} & True & 26.4 \\
19 & \texttt{Killing spinors of S\textasciicircum{}4 (radius r0, frame e\textasciicircum{}a = (r0, r0 sin x4, ...) dx\textasciicircum{}(a+3), gt\_\allowbreak{}a = I rho\_\allowbreak{}(a+3)): eps\_\allowbreak{}s = exp(s (i/\allowbreak{}2) x4 gt\_\allowbreak{}1) exp((1/\allowbreak{}2) x5 gt\_\allowbreak{}1 gt\_\allowbreak{}2) exp((1/\allowbreak{}2) x6 gt\_\allowbreak{}2 gt\_\allowbreak{}3) exp((1/\allowbreak{}2) x7 gt\_\allowbreak{}3 gt\_\allowbreak{}4) eps0 satisfies nabla\_\allowbreak{}mu eps\_\allowbreak{}s == s (i/\allowbreak{}(2 r0)) e\_\allowbreak{}mu\textasciicircum{}a gt\_\allowbreak{}a eps\_\allowbreak{}s (s = +-1, all eps0), so the Dirac operator gives (2 s i/\allowbreak{}r0) eps\_\allowbreak{}s: |lambda| = 2/\allowbreak{}r0, Friedrich's bound n R/\allowbreak{}(4(n - 1)) = 4/\allowbreak{}r0\textasciicircum{}2} & True & 26.4 \\
20 & \texttt{Dirac spectrum on S\textasciicircum{}4: Weyl's law for spinors (4 components) N(|lambda| <= L/\allowbreak{}r) \textasciitilde{} 4 V (L/\allowbreak{}r)\textasciicircum{}4/\allowbreak{}(32 Pi\textasciicircum{}2) = L\textasciicircum{}4/\allowbreak{}3 agrees with the degeneracies 2 x 4 Binomial[l + 3, 3] of the eigenvalues +-i(l + 2)/\allowbreak{}r (to 2.0 10\textasciicircum{}-3 at L = 1000)} & True & 26.4 \\
21 & \texttt{the sphere with the volume of the torus today (V = M\_\allowbreak{}Pl\textasciicircum{}2/\allowbreak{}M8\textasciicircum{}6): r = (3 V/\allowbreak{}(8 Pi\textasciicircum{}2))\textasciicircum{}(1/\allowbreak{}4) = 688.9 GeV\textasciicircum{}-1 = 1.359 10\textasciicircum{}-13 m; lowest scalar levels 4.03 MeV (torus, 8 states) and 2.90 MeV (sphere, l = 1, 5 states); lowest dirac16complex level on the sphere 2/\allowbreak{}r = 2.90 MeV (8 Dirac fermions; the torus has 4 massless ones for m = 0); levels below 5 GeV: torus 1.17335 10\textasciicircum{}13, sphere 1.17375 10\textasciicircum{}13, Weyl 1.17335 10\textasciicircum{}13} & True & 26.4 \\
22 & \texttt{RS1 in five dimensions (mostly minus): G\textasciicircum{}mu\_\allowbreak{}nu = 3 A'' - 6 A'\textasciicircum{}2, G\textasciicircum{}y\_\allowbreak{}y = -6 A'\textasciicircum{}2; bulk A = k y needs Lambda = -6 k\textasciicircum{}2; the delta functions A'' = +-2 k delta at the branes need the tensions f\_\allowbreak{}UV = +6 k/\allowbreak{}kappa, f\_\allowbreak{}IR = -6 k/\allowbreak{}kappa (UV positive)} & True & 26.5 \\
23 & \texttt{Israel junction at the UV brane of RS1: with K\textasciicircum{}mu\_\allowbreak{}nu = (1/\allowbreak{}2) g\textasciicircum{}mu alpha d\_\allowbreak{}y g\_\allowbreak{}alpha nu = -A' delta\textasciicircum{}mu\_\allowbreak{}nu (normal d\_\allowbreak{}y), [K\textasciicircum{}mu\_\allowbreak{}nu] - delta\textasciicircum{}mu\_\allowbreak{}nu [K] = 6 k = (coefficient of A'' in G\textasciicircum{}mu\_\allowbreak{}nu) x (jump of A'), the delta-function part of G\textasciicircum{}mu\_\allowbreak{}nu; in the signature (-,+,+,+,+) of the note G\textasciicircum{}A\_\allowbreak{}B changes sign, so there [K\_\allowbreak{}mu nu - g\_\allowbreak{}mu nu K] = -kappa S\_\allowbreak{}mu nu (the note uses +kappa S\_\allowbreak{}mu nu, hence its reversed tensions)} & True & 26.5 \\
24 & \texttt{RS1: M\_\allowbreak{}Pl\textasciicircum{}2 = (1/\allowbreak{}kappa) Int\_\allowbreak{}(-Pi R)\textasciicircum{}(Pi R) e\textasciicircum{}(-2 k |y|) dy = (1 - e\textasciicircum{}(-2 Pi k R))/\allowbreak{}(kappa k)} & True & 26.5 \\
25 & \texttt{RS1 radion: with R promoted to T(x), the bulk action is (3/\allowbreak{}(kappa k)) (d e\textasciicircum{}(-k Pi T))\textasciicircum{}2 + V(T) + (a total derivative), so the canonical radion is mu = Sqrt[6/\allowbreak{}(kappa k)] e\textasciicircum{}(-k Pi T) (the note: Sqrt[12/\allowbreak{}(kappa k)])} & True & 26.5 \\
26 & \texttt{eight dimensions, warped along one extra direction (ds\textasciicircum{}2 = e\textasciicircum{}(-2 A(y)) eta - dy\textasciicircum{}2 - dz.dz): G\textasciicircum{}mu\_\allowbreak{}nu = 3 A'' - 6 A'\textasciicircum{}2, G\textasciicircum{}y\_\allowbreak{}y = -6 A'\textasciicircum{}2, G\textasciicircum{}z\_\allowbreak{}z = 4 A'' - 10 A'\textasciicircum{}2; with only Lambda (all G\textasciicircum{}A\_\allowbreak{}A equal) A'' = 0 and A' = 0: the note's five-dimensional equations do not carry over} & True & 26.5 \\
27 & \texttt{Goldberger-Wise in RS1: the static bulk scalar (Sqrt|G| = e\textasciicircum{}(-4 k y), G\textasciicircum{}yy = -1) obeys Phi'' - 4 k Phi' - m\textasciicircum{}2 Phi = 0, solved by e\textasciicircum{}(2 k y)(C1 e\textasciicircum{}(nu k y) + C2 e\textasciicircum{}(-nu k y)), nu = Sqrt[4 + m\textasciicircum{}2/\allowbreak{}k\textasciicircum{}2] = 2 + eps + O(m\textasciicircum{}4), eps = m\textasciicircum{}2/\allowbreak{}(4 k\textasciicircum{}2)} & True & 26.5 \\
28 & \texttt{Goldberger-Wise radion mass: V(mu) = 4 k (mu/\allowbreak{}f)\textasciicircum{}4 (v\_\allowbreak{}B - v\_\allowbreak{}A (mu/\allowbreak{}f)\textasciicircum{}eps)\textasciicircum{}2 has V'(<mu>) = 0 at (<mu>/\allowbreak{}f)\textasciicircum{}eps = v\_\allowbreak{}B/\allowbreak{}v\_\allowbreak{}A (k Pi <T> = (4 k\textasciicircum{}2/\allowbreak{}m\textasciicircum{}2) Log[v\_\allowbreak{}A/\allowbreak{}v\_\allowbreak{}B]) and V''(<mu>) = 8 k eps\textasciicircum{}2 v\_\allowbreak{}B\textasciicircum{}2 <mu>\textasciicircum{}2/\allowbreak{}f\textasciicircum{}4; with f\textasciicircum{}2 = 6/\allowbreak{}(kappa k): m\_\allowbreak{}r\textasciicircum{}2 = (4/\allowbreak{}3) kappa k\textasciicircum{}2 eps\textasciicircum{}2 v\_\allowbreak{}B\textasciicircum{}2 (<mu>/\allowbreak{}f)\textasciicircum{}2 (the note's f\textasciicircum{}2 = 12/\allowbreak{}(kappa k) gives 2/\allowbreak{}3)} & True & 26.5 \\
29 & \texttt{conformally flat brane metric h = Omega\textasciicircum{}2 eta (four dimensions, Omega(x0, x1)): Sqrt|h| R(h) == -6 Omega box Omega with box = d0\textasciicircum{}2 - d1\textasciicircum{}2 (curvature conventions of 7+1.nb)} & True & 26.6 \\
30 & \texttt{the Higgs boson on the brane (h = Omega\textasciicircum{}2 eta, Omega = 1 + r/\allowbreak{}Lambda + a2 r\textasciicircum{}2, phi\_\allowbreak{}H = v + h, gamma = v/\allowbreak{}Lambda): the terms linear in (h, r) are total derivatives, and to second order L == (1/\allowbreak{}2)(dh)\textasciicircum{}2 + (1/\allowbreak{}2)(1 + 6 xi gamma\textasciicircum{}2)(dr)\textasciicircum{}2 + 6 xi gamma dh.dr - (1/\allowbreak{}2) m\_\allowbreak{}h0\textasciicircum{}2 h\textasciicircum{}2 - (1/\allowbreak{}2) m\_\allowbreak{}r0\textasciicircum{}2 r\textasciicircum{}2 up to total derivatives (Euler-Lagrange expressions in x0 and x1); a2 drops out} & True & 26.6 \\
31 & \texttt{canonical kinetic terms: Transpose[S].K.S == 1 for S = \{\{1, -6 xi gamma/\allowbreak{}Z\}, \{0, 1/\allowbreak{}Z\}\} if and only if Z\textasciicircum{}2 = 1 + 6 xi gamma\textasciicircum{}2 (1 - 6 xi) (the note's Z); with the note's literal kinetic terms (no 6 xi gamma\textasciicircum{}2 for the radion) it would be Z\textasciicircum{}2 = 1 - 36 xi\textasciicircum{}2 gamma\textasciicircum{}2} & True & 26.6 \\
32 & \texttt{mass mixing: with (h', r') = Q(theta).(h\_\allowbreak{}m, r\_\allowbreak{}m), Q = \{\{cos, -sin\}, \{sin, cos\}\}, the off-diagonal element of Transpose[Q].M'.Q is -(1/\allowbreak{}2)(M'\_\allowbreak{}11 - M'\_\allowbreak{}22) sin 2 theta + M'\_\allowbreak{}12 cos 2 theta, zero for tan 2 theta = 2 M'\_\allowbreak{}12/\allowbreak{}(M'\_\allowbreak{}11 - M'\_\allowbreak{}22) = 12 gamma xi Z m\_\allowbreak{}h0\textasciicircum{}2/\allowbreak{}(m\_\allowbreak{}r0\textasciicircum{}2 - m\_\allowbreak{}h0\textasciicircum{}2 (Z\textasciicircum{}2 - 36 xi\textasciicircum{}2 gamma\textasciicircum{}2)) (the note's formula); with the inverse matrix (the note's display) it is +(1/\allowbreak{}2)(M'\_\allowbreak{}11 - M'\_\allowbreak{}22) sin 2 theta + M'\_\allowbreak{}12 cos 2 theta (theta -> -theta); (h, r) = S.Q.(h\_\allowbreak{}m, r\_\allowbreak{}m) gives C\_\allowbreak{}h = cos theta - (6 xi gamma/\allowbreak{}Z) sin theta, C\_\allowbreak{}r = sin theta/\allowbreak{}Z} & True & 26.6 \\
33 & \texttt{Higgs-radion mixing in 7+1 (gamma = v/\allowbreak{}(Sqrt[3] M\_\allowbreak{}Pl), xi = 1, m\_\allowbreak{}r0 = 10 GeV, m\_\allowbreak{}h0 = 125.20 GeV; 60-digit arithmetic): |Z - 1| < 10\textasciicircum{}-30, |theta| = 3.5 10\textasciicircum{}-16, the rates of h -> gamma gamma and h -> g g change by less than 10\textasciicircum{}-30; in RS1 with <mu> = 3 TeV gamma = v/\allowbreak{}(3 TeV) = 0.082} & True & 26.6 \\*
\end{longtable}}
The Python program’s boolean checks:

{\footnotesize
\begin{longtable}{>{\raggedright\arraybackslash}p{\dimexpr 0.0500\linewidth-2\tabcolsep\relax}>{\raggedright\arraybackslash}p{\dimexpr 0.8300\linewidth-2\tabcolsep\relax}>{\raggedright\arraybackslash}p{\dimexpr 0.1200\linewidth-2\tabcolsep\relax}}
\toprule
\textbf{\#} & \textbf{check} & \textbf{result} \\
\midrule\endfirsthead
\toprule
\textbf{\#} & \textbf{check} & \textbf{result} \\
\midrule\endhead
\endfoot
\bottomrule\endlastfoot
1 & \texttt{the 4-sphere of radius r0: Ricci\_\allowbreak{}ij = 3 gamma\_\allowbreak{}ij, off-diagonal 0, R = -12/\allowbreak{}r0\textasciicircum{}2 in the signature of 7+1.nb; volume of the unit S\textasciicircum{}4 = 8 pi\textasciicircum{}2/\allowbreak{}3} & True \\
2 & \texttt{radion on T\textasciicircum{}4: the reduced action equals -(M\_\allowbreak{}Pl\textasciicircum{}2/\allowbreak{}2) R\_\allowbreak{}E + 6 M\_\allowbreak{}Pl\textasciicircum{}2 (d beta)\textasciicircum{}2 - e\textasciicircum{}(-4 beta)(M\_\allowbreak{}Pl\textasciicircum{}2 Lambda + V0 U) up to a total derivative (Euler-Lagrange for N, a, beta), M\_\allowbreak{}Pl\textasciicircum{}2 = V0/\allowbreak{}kappa} & True \\
3 & \texttt{radion on S\textasciicircum{}4: the same with the potential -(M\_\allowbreak{}Pl\textasciicircum{}2/\allowbreak{}2) e\textasciicircum{}(-4 beta) 12/\allowbreak{}(r0\textasciicircum{}2 e\textasciicircum{}(2 beta)) from the curvature of the sphere, V0 = (8 pi\textasciicircum{}2/\allowbreak{}3) r0\textasciicircum{}4} & True \\
4 & \texttt{n extra dimensions (T\textasciicircum{}n, n = 1..4): kinetic term (M\_\allowbreak{}Pl\textasciicircum{}2/\allowbreak{}2)(n(n+2)/\allowbreak{}2)(dB)\textasciicircum{}2; the note's beta\_\allowbreak{}n = sqrt(2/\allowbreak{}(n(n+2))) gives (M\_\allowbreak{}Pl\textasciicircum{}2/\allowbreak{}2)(d phi)\textasciicircum{}2, not (M\_\allowbreak{}Pl\textasciicircum{}2/\allowbreak{}2)(1/\allowbreak{}2)(d phi)\textasciicircum{}2; beta\_\allowbreak{}n = 1/\allowbreak{}sqrt(n(n+2)) gives the 1/\allowbreak{}2} & True \\
5 & \texttt{the note's eight-dimensional alpha = -1/\allowbreak{}sqrt(3), beta = 1/\allowbreak{}sqrt(12): 2 alpha + 4 beta = 0, 4 beta = 2/\allowbreak{}sqrt(3), sqrt(12) beta = 1} & True \\
6 & \texttt{the kinetic coefficient solved from the reduction (n = 1..4) is n(n+2)/\allowbreak{}2} & True \\
7 & \texttt{the cosmological metric with the sphere of radius b: G\textasciicircum{}A\_\allowbreak{}A = torus + (6, 6, 6, 6, 3, 3, 3, 3)/\allowbreak{}b\textasciicircum{}2, R = torus - 12/\allowbreak{}b\textasciicircum{}2} & True \\
8 & \texttt{radion equation on S\textasciicircum{}4: Hb' + 3 Ha Hb + 4 Hb\textasciicircum{}2 = kappa (eps/\allowbreak{}6 - pa/\allowbreak{}2 + pb/\allowbreak{}3) + Lambda/\allowbreak{}3 - 3/\allowbreak{}b\textasciicircum{}2} & True \\
9 & \texttt{radion on T\textasciicircum{}4 with U: flat vacuum Lambda + kappa U0 = 0, U1 = 0; m\_\allowbreak{}r\textasciicircum{}2 = kappa K\_\allowbreak{}U/\allowbreak{}12} & True \\
10 & \texttt{radion on S\textasciicircum{}4 with U: Lambda + kappa U0 = 6/\allowbreak{}r0\textasciicircum{}2, kappa U1 = -12/\allowbreak{}r0\textasciicircum{}2, m\_\allowbreak{}r\textasciicircum{}2 = kappa K\_\allowbreak{}U/\allowbreak{}12 - 2/\allowbreak{}r0\textasciicircum{}2, from V and from the 8D radion equation} & True \\
11 & \texttt{S\textasciicircum{}4 with Lambda only: dS4 x S\textasciicircum{}4 (r0\textasciicircum{}2 = 9/\allowbreak{}Lambda, H\textasciicircum{}2 = 1/\allowbreak{}r0\textasciicircum{}2) solves the 00, 11, 44 equations; maximum of V with m\textasciicircum{}2 = -6 H\textasciicircum{}2, also from the 8D radion equation} & True \\
12 & \texttt{S\textasciicircum{}4 with a 4-form flux: eps = rho\_\allowbreak{}F, pa = -rho\_\allowbreak{}F, pb = +rho\_\allowbreak{}F; Minkowski needs Lambda = 9/\allowbreak{}(2 r0\textasciicircum{}2), kappa rho\_\allowbreak{}F = 3/\allowbreak{}(2 r0\textasciicircum{}2); m\textasciicircum{}2 = 6/\allowbreak{}r0\textasciicircum{}2 from V and from the 8D radion equation} & True \\
13 & \texttt{the note's Goldberger-Wise potential: for a static Phi(y) the Einstein-frame energy carries e\textasciicircum{}(-4 beta) (gradient part e\textasciicircum{}(-6 beta)), not e\textasciicircum{}(+4 beta)} & True \\
14 & \texttt{radion coupling to a mass on the brane: -M e\textasciicircum{}(-n beta/\allowbreak{}2) = -M + (n/\allowbreak{}2) M beta; coupling sqrt(n/\allowbreak{}(2(n+2)))/\allowbreak{}M\_\allowbreak{}Pl = 1/\allowbreak{}(sqrt(3) M\_\allowbreak{}Pl) for n = 4; fifth force 2 (coupling M\_\allowbreak{}Pl)\textasciicircum{}2 = n/\allowbreak{}(n+2) = 2/\allowbreak{}3} & True \\
15 & \texttt{Eot-Wash: m\_\allowbreak{}r > hbar c/\allowbreak{}38.6 micrometres = 5.1 meV; range of a 10 GeV radion 2.0e-17 m; gamma = v/\allowbreak{}(sqrt(3) M\_\allowbreak{}Pl) = 5.8e-17} & True \\
16 & \texttt{zonal harmonics on S\textasciicircum{}4 (l = 0..6): eigenvalue -l(l+3), C\_\allowbreak{}l(1) = (l+1)(l+2)/\allowbreak{}2, Int (C\_\allowbreak{}l/\allowbreak{}C\_\allowbreak{}l(1))\textasciicircum{}2 = V/\allowbreak{}d\_\allowbreak{}l; d\_\allowbreak{}l = binomial(l+4,4) - binomial(l+2,4) (l <= 20)} & True \\
17 & \texttt{Weyl's law on S\textasciicircum{}4: N(x)/\allowbreak{}(x\textasciicircum{}4/\allowbreak{}12) - 1 = -1/\allowbreak{}x\textasciicircum{}2 at x = 10, 100, 1000} & True \\
18 & \texttt{the brane at the north pole of S\textasciicircum{}4: the embedding induces r0\textasciicircum{}2 gamma; north pole (0,0,0,0,r0); 6 of the 10 rotations of R\textasciicircum{}5 fix it (SO(4)), 4 move it} & True \\
19 & \texttt{Killing spinors of S\textasciicircum{}4: nabla eps\_\allowbreak{}s = s (i/\allowbreak{}(2 r0)) gt eps\_\allowbreak{}s and D eps\_\allowbreak{}s = (2 s i/\allowbreak{}r0) eps\_\allowbreak{}s (s = +-1; residual < 1e-25 at four random points, 30 digits); Friedrich bound 4/\allowbreak{}r0\textasciicircum{}2} & True \\
20 & \texttt{Dirac spectrum on S\textasciicircum{}4: Weyl's law for spinors L\textasciicircum{}4/\allowbreak{}3 versus the degeneracies 2 x 4 binomial(l+3,3) (2.0e-3 at L = 1000)} & True \\
21 & \texttt{the sphere with the volume of the torus today: r = 688.9 GeV\textasciicircum{}-1; lowest levels 4.03 MeV (torus) and 2.90 MeV (sphere); dirac16complex on the sphere 2.90 MeV; levels below 5 GeV agree with Weyl} & True \\
22 & \texttt{RS1 (5D, mostly minus): G\textasciicircum{}mu\_\allowbreak{}nu = 3 A'' - 6 A'\textasciicircum{}2, G\textasciicircum{}y\_\allowbreak{}y = -6 A'\textasciicircum{}2; bulk Lambda = -6 k\textasciicircum{}2; tensions f\_\allowbreak{}UV = +6k/\allowbreak{}kappa, f\_\allowbreak{}IR = -6k/\allowbreak{}kappa} & True \\
23 & \texttt{Israel junction at the UV brane: K\textasciicircum{}mu\_\allowbreak{}nu = -A' delta, [K\textasciicircum{}mu\_\allowbreak{}nu] - delta [K] = 6k = 3 x (jump of A'), the delta part of G\textasciicircum{}mu\_\allowbreak{}nu; in the signature (-,+,+,+,+) [K - g K] = -kappa S} & True \\
24 & \texttt{RS1: M\_\allowbreak{}Pl\textasciicircum{}2 = (1/\allowbreak{}kappa) Int e\textasciicircum{}(-2k|y|) dy over [-pi R, pi R] = (1 - e\textasciicircum{}(-2 pi k R))/\allowbreak{}(kappa k)} & True \\
25 & \texttt{RS1 radion: kinetic term (3/\allowbreak{}(kappa k)) (d e\textasciicircum{}(-k pi T))\textasciicircum{}2, canonical mu = sqrt(6/\allowbreak{}(kappa k)) e\textasciicircum{}(-k pi T)} & True \\
26 & \texttt{8D warped along one extra direction: G\textasciicircum{}y\_\allowbreak{}y = -6 A'\textasciicircum{}2, G\textasciicircum{}z\_\allowbreak{}z = 4 A'' - 10 A'\textasciicircum{}2; with only Lambda A'' = 0 and A' = 0} & True \\
27 & \texttt{Goldberger-Wise: Phi'' - 4k Phi' - m\textasciicircum{}2 Phi = 0 from the bulk action, solved by e\textasciicircum{}(2ky)(C1 e\textasciicircum{}(nu k y) + C2 e\textasciicircum{}(-nu k y)), nu = sqrt(4 + m\textasciicircum{}2/\allowbreak{}k\textasciicircum{}2) = 2 + eps; V'(<mu>) = 0 at (<mu>/\allowbreak{}f)\textasciicircum{}eps = v\_\allowbreak{}B/\allowbreak{}v\_\allowbreak{}A, V''(<mu>) = 8 k eps\textasciicircum{}2 v\_\allowbreak{}B\textasciicircum{}2 <mu>\textasciicircum{}2/\allowbreak{}f\textasciicircum{}4; m\_\allowbreak{}r\textasciicircum{}2 = (4/\allowbreak{}3) kappa k\textasciicircum{}2 eps\textasciicircum{}2 v\_\allowbreak{}B\textasciicircum{}2 (<mu>/\allowbreak{}f)\textasciicircum{}2 (note: 2/\allowbreak{}3)} & True \\
28 & \texttt{conformally flat h = Omega\textasciicircum{}2 eta: sqrt|h| R(h) = -6 Omega box Omega, box = d0\textasciicircum{}2 - d1\textasciicircum{}2} & True \\
29 & \texttt{the Higgs boson on the brane: linear terms are total derivatives; to second order L = (1/\allowbreak{}2)(dh)\textasciicircum{}2 + (1/\allowbreak{}2)(1 + 6 xi gamma\textasciicircum{}2)(dr)\textasciicircum{}2 + 6 xi gamma dh.dr - (1/\allowbreak{}2) m\_\allowbreak{}h0\textasciicircum{}2 h\textasciicircum{}2 - (1/\allowbreak{}2) m\_\allowbreak{}r0\textasciicircum{}2 r\textasciicircum{}2 up to total derivatives; o\_\allowbreak{}2 drops out} & True \\
30 & \texttt{the shift h = h' - (6 xi gamma/\allowbreak{}Z) r', r = r'/\allowbreak{}Z is canonical iff Z\textasciicircum{}2 = 1 + 6 xi gamma\textasciicircum{}2 (1 - 6 xi) (note's literal kinetic terms: 1 - 36 xi\textasciicircum{}2 gamma\textasciicircum{}2); tan 2 theta = 12 gamma xi Z m\_\allowbreak{}h0\textasciicircum{}2/\allowbreak{}(m\_\allowbreak{}r0\textasciicircum{}2 - m\_\allowbreak{}h0\textasciicircum{}2 (Z\textasciicircum{}2 - 36 xi\textasciicircum{}2 gamma\textasciicircum{}2)) diagonalizes with (h', r') = Q(theta)(h\_\allowbreak{}m, r\_\allowbreak{}m); the inverse matrix needs -theta; C\_\allowbreak{}h = cos - (6 xi gamma/\allowbreak{}Z) sin, C\_\allowbreak{}r = sin/\allowbreak{}Z} & True \\
31 & \texttt{Higgs-radion mixing in 7+1 (gamma = v/\allowbreak{}(sqrt(3) M\_\allowbreak{}Pl), xi = 1, m\_\allowbreak{}r0 = 10 GeV): |Z - 1| < 1e-30, |theta| = 3.5e-16, rates change by < 1e-30; warped gamma = v/\allowbreak{}(3 TeV) = 0.082} & True \\*
\end{longtable}}
The numbers printed by both programs (largest relative difference per group; integers must agree exactly):

{\footnotesize
\begin{longtable}{>{\raggedright\arraybackslash}p{\dimexpr 0.5200\linewidth-2\tabcolsep\relax}>{\raggedright\arraybackslash}p{\dimexpr 0.1000\linewidth-2\tabcolsep\relax}>{\raggedright\arraybackslash}p{\dimexpr 0.2200\linewidth-2\tabcolsep\relax}>{\raggedright\arraybackslash}p{\dimexpr 0.1600\linewidth-2\tabcolsep\relax}}
\toprule
\textbf{quantity} & \textbf{values} & \textbf{max. relative difference} & \textbf{tolerance} \\
\midrule\endfirsthead
\toprule
\textbf{quantity} & \textbf{values} & \textbf{max. relative difference} & \textbf{tolerance} \\
\midrule\endhead
\endfoot
\bottomrule\endlastfoot
kinetic coefficient $n(n+2)/2$, $n=1,\dots,4$ (solved from the reduction) & 4 & 0 & exact \\
$\beta_n^2$ of the note and the canonical $\beta_n^2$, $n=4$ & 2 & 0 & $1\times10^{-15}$ \\
$S^4$ with $U$: $(\Lambda+\kappa U_0)r_0^2$, $\kappa U_1r_0^2$, $(m_r^2-\kappa K_U/12)r_0^2$ & 3 & 0 & exact \\
$S^4$ with $\Lambda$ only: $\Lambda r_0^2$, $H^2r_0^2$, $m^2/H^2$ & 3 & 0 & exact \\
$S^4$ with a 4-form flux: $\Lambda r_0^2$, $\kappa\rho_Fr_0^2$, $m^2r_0^2$ & 3 & 0 & $1\times10^{-15}$ \\
E\textbackslash{}"ot-Wash bound, range of a 10 GeV radion, $\gamma=v/(\sqrt3M_{\rm Pl})$ & 3 & $2.1\times10^{-16}$ & $1\times10^{-10}$ \\
multiplicities $d_l$ on $S^4$, $l=0,\dots,8$ & 9 & 0 & exact \\
Weyl's law on $S^4$: $N/(x^4/12)-1$ at $x=10,100,1000$ & 3 & $8.2\times10^{-11}$ & $1\times10^{-9}$ \\
Dirac degeneracies on $S^4$ and the spinor Weyl ratio & 7 & $7.0\times10^{-14}$ & $1\times10^{-9}$ \\
today: $r_{S^4}$, lowest masses, numbers of levels below 5 GeV & 8 & $2.2\times10^{-16}$ & $1\times10^{-12}$ \\
RS1: $3$, $f_{\rm UV}\kappa/k$, $f_{\rm IR}\kappa/k$, $f^2\kappa k$ & 4 & 0 & exact \\
strength of the warped coupling over the flat one & 1 & $1.8\times10^{-16}$ & $1\times10^{-12}$ \\
Goldberger--Wise mass coefficient: $\frac43$ and the note's $\frac23$ & 2 & 0 & $1\times10^{-15}$ \\
Higgs--radion: $c=-6$, $K_{hr}/(\xi\gamma)$, $(K_{rr}-1)/(\xi\gamma^2)$ & 3 & 0 & exact \\
Higgs--radion mixing in 7+1: $Z-1$, $\theta$, rate changes; warped $\gamma$ & 5 & $4.4\times10^{-16}$ & $1\times10^{-12}$ \\*
\end{longtable}}
The results quoted from the other sections of 7+1.nb rest on these checks of the notebook, all evaluated in the same run:

{\footnotesize
\begin{longtable}{>{\raggedright\arraybackslash}p{\dimexpr 0.0500\linewidth-2\tabcolsep\relax}>{\raggedright\arraybackslash}p{\dimexpr 0.7700\linewidth-2\tabcolsep\relax}>{\raggedright\arraybackslash}p{\dimexpr 0.0800\linewidth-2\tabcolsep\relax}>{\raggedright\arraybackslash}p{\dimexpr 0.1000\linewidth-2\tabcolsep\relax}}
\toprule
\textbf{\#} & \textbf{check} & \textbf{result} & \textbf{section} \\
\midrule\endfirsthead
\toprule
\textbf{\#} & \textbf{check} & \textbf{result} & \textbf{section} \\
\midrule\endhead
\endfoot
\bottomrule\endlastfoot
1 & \texttt{stabilizing potential: W = a\textasciicircum{}3 b\textasciicircum{}4 U(b) gives eps = U, pa = -U, pb = -U - (b/\allowbreak{}4) U'(b)} & True & 20.2 \\
2 & \texttt{radion equation of Einstein gravity: Hb' + 3 Ha Hb + 4 Hb\textasciicircum{}2 == kappa (eps/\allowbreak{}6 - pa/\allowbreak{}2 + pb/\allowbreak{}3) + Lambda/\allowbreak{}3 (constraint used)} & True & 20.3 \\
3 & \texttt{radion mass: U = U0 + (K\_\allowbreak{}U/\allowbreak{}2) beta\textasciicircum{}2, b U' = K\_\allowbreak{}U beta gives the restoring term -(kappa K\_\allowbreak{}U/\allowbreak{}12) beta, m\_\allowbreak{}r\textasciicircum{}2 = kappa K\_\allowbreak{}U/\allowbreak{}12} & True & 20.3 \\
4 & \texttt{static auxiliary space: (kappa/\allowbreak{}4) b U' == Lambda + kappa U + 2 kappa f\_\allowbreak{}W/\allowbreak{}V\_\allowbreak{}perp + (kappa/\allowbreak{}2)(eps - 3pa) + kappa pb - 12 alpha2 Ha\textasciicircum{}2 (Ha' + Ha\textasciicircum{}2)} & True & 20.3 \\
5 & \texttt{Einstein-frame form of the constraint: 3 (Ha + 2 Hb)\textasciicircum{}2 == 6 Hb\textasciicircum{}2 + (E\textasciicircum{}(1)0\_\allowbreak{}0 of section 12.6), i.e. 3 M\_\allowbreak{}Pl\textasciicircum{}2 H\_\allowbreak{}E\textasciicircum{}2 == phi'\textasciicircum{}2/\allowbreak{}2 + rho with phi = Sqrt[12] M\_\allowbreak{}Pl beta} & True & 20.3 \\
6 & \texttt{Kasner constants: P\_\allowbreak{}E(u1) == 0, p1 == (1 + 2 Sqrt[2])/\allowbreak{}7, q1 == (2 - 3 Sqrt[2])/\allowbreak{}14; u2 in (-1/\allowbreak{}3, 0) with P\_\allowbreak{}G(u2) == 0, 3 p2 + 4 q2 == 3} & True & 21.1 \\
7 & \texttt{without the stabilizing potential the gas-dominated Einstein regime decompactifies: tau Ha -> 0, tau Hb -> 2/\allowbreak{}5, pb/\allowbreak{}eps -> 1/\allowbreak{}4 (Ha = 10\textasciicircum{}-14)} & True & 21.3 \\
8 & \texttt{capture: the radion oscillates like a four-dimensional massive field in a matter-dominated universe: with tau\_\allowbreak{}E from <H\_\allowbreak{}E> = 2/\allowbreak{}(3 (tau - tau\_\allowbreak{}E)) over the last period, beta = A sin(m\_\allowbreak{}r tau + vartheta)/\allowbreak{}(m\_\allowbreak{}r (tau - tau\_\allowbreak{}E)) to 2 10\textasciicircum{}-2, <H\_\allowbreak{}E>(tau - tau\_\allowbreak{}E) = 2/\allowbreak{}3 to 10\textasciicircum{}-2 at tau = 10\textasciicircum{}5 and 2 10\textasciicircum{}5, the gas share constant to 2 10\textasciicircum{}-2} & True & 22.1 \\
9 & \texttt{through brane fields alone a 10 GeV radion would decay far too late: T\_\allowbreak{}RH < 10\textasciicircum{}-6 GeV (nucleosynthesis needs 4 10\textasciicircum{}-3 GeV); 10\textasciicircum{}13 Kaluza-Klein channels below m\_\allowbreak{}r/\allowbreak{}2} & True & 22.2 \\
10 & \texttt{the radion decays into the bulk: the lattice sum over the Kaluza-Klein modes of dirac16complex (m << m\_\allowbreak{}r/\allowbreak{}2) equals its continuum limit to 10\textasciicircum{}-6, the number of momenta to 10\textasciicircum{}-5, and Gamma\_\allowbreak{}bulk/\allowbreak{}Gamma\_\allowbreak{}brane > 10\textasciicircum{}12} & True & 22.2 \\
11 & \texttt{radion equilibrium in de Sitter: 3 m\_\allowbreak{}r\textasciicircum{}2 beta\_\allowbreak{}eq == 3 H\textasciicircum{}2 - 12 alpha2 H\textasciicircum{}4 + O(beta\_\allowbreak{}eq\textasciicircum{}2)} & True & 23.2 \\
12 & \texttt{eight-dimensional de Sitter (no stabilizer): all three equations read 21 H\textasciicircum{}2 - 420 alpha2 H\textasciicircum{}4 + 2520 alpha3 H\textasciicircum{}6 == Lambda} & True & 23.2 \\
13 & \texttt{ADD relation: kappa8 rho4/\allowbreak{}V\_\allowbreak{}perp == 8 Pi G4 rho4 with G4 = G8/\allowbreak{}V\_\allowbreak{}perp (kappa8 = 8 Pi G8)} & True & 17.5 \\
14 & \texttt{the auxiliary torus today (M8 = 10 TeV, M\_\allowbreak{}Pl = 2.435 10\textasciicircum{}18 GeV): V\_\allowbreak{}perp = M\_\allowbreak{}Pl\textasciicircum{}2/\allowbreak{}M8\textasciicircum{}6 = 5.929 10\textasciicircum{}12 GeV\textasciicircum{}-4, l\_\allowbreak{}perp = V\_\allowbreak{}perp\textasciicircum{}(1/\allowbreak{}4) = b0 Lc = 1560.4 GeV\textasciicircum{}-1 = 3.079 10\textasciicircum{}-13 m, lowest Kaluza-Klein mass 2 Pi/\allowbreak{}l\_\allowbreak{}perp = 4.027 MeV; N(mu) = V\_\allowbreak{}perp mu\textasciicircum{}4/\allowbreak{}(32 Pi\textasciicircum{}2) = 1.17 10\textasciicircum{}13 levels below m\_\allowbreak{}r/\allowbreak{}2 = 5 GeV (the lattice count of section 22.2 agrees to 10\textasciicircum{}-5) and N(M8) = M\_\allowbreak{}Pl\textasciicircum{}2/\allowbreak{}(32 Pi\textasciicircum{}2 M8\textasciicircum{}2)} & True & 25.3 \\
15 & \texttt{graviton + radion: far from the source the images add up to (5 kappa M/\allowbreak{}(32 Pi\textasciicircum{}3)) (4 Pi\textasciicircum{}2/\allowbreak{}(3 V\_\allowbreak{}perp r)) == (5/\allowbreak{}3) G4 M/\allowbreak{}r with G4 = kappa/\allowbreak{}(8 Pi V\_\allowbreak{}perp); the radion of section 20.3, coupled to matter by (phi/\allowbreak{}(Sqrt[3] M\_\allowbreak{}Pl)) T, gives the potential -(1/\allowbreak{}(Sqrt[3] M\_\allowbreak{}Pl))\textasciicircum{}2 M M' Exp[-m\_\allowbreak{}r r]/\allowbreak{}(4 Pi r), and (1/\allowbreak{}(Sqrt[3] M\_\allowbreak{}Pl))\textasciicircum{}2/\allowbreak{}(4 Pi) == (2/\allowbreak{}3) G4 (8 Pi G4 = 1/\allowbreak{}M\_\allowbreak{}Pl\textasciicircum{}2): graviton 1 + massless radion 2/\allowbreak{}3 = 5/\allowbreak{}3} & True & 25.4 \\*
\end{longtable}}
\Needspace{12\baselineskip}\section*{14. How to reproduce}\addcontentsline{toc}{section}{14. How to reproduce}
\begin{enumerate}
\item Write the unevaluated notebook with \texttt{python3 tools/\allowbreak{}make\_\allowbreak{}7p1\_\allowbreak{}nb.py}; section 26 is generated by \texttt{tools/\allowbreak{}nb\_\allowbreak{}part\_\allowbreak{}radion.py}.
\item Evaluate 7+1.nb in a fresh kernel of Wolfram 15.0.1 (Evaluation \texorpdfstring{\ensuremath{\triangleright}}{▸} Quit Kernel, then Evaluation \texorpdfstring{\ensuremath{\triangleright}}{▸} Evaluate Notebook; or \texttt{Get["D:/\allowbreak{}Developer/\allowbreak{}github/\allowbreak{}7+1/\allowbreak{}tools/\allowbreak{}run\_\allowbreak{}eval.wl"]}) and save it.
\item Run \texttt{python3 tools/\allowbreak{}ref\_\allowbreak{}radion.py}: it repeats the checks independently (exit code 0 if all pass).
\item Run \texttt{python3 tools/\allowbreak{}make\_\allowbreak{}docs.py}: it regenerates this document, its symbol dictionary and its provenance record (and all other documents).
\end{enumerate}
\Needspace{12\baselineskip}\section*{15. References}\addcontentsline{toc}{section}{15. References}
\begin{itemize}
\item Th. Kaluza, Sitzungsber. Preuss. Akad. Wiss. (1921) 966; O. Klein, Z. Phys. 37 (1926) 895 — compact extra dimensions and the size modulus.
\item C. N. Pope, \emph{Kaluza–Klein Theory}, lecture notes (Texas A\&M) — the reduction $e^{2\alpha\varphi}g+e^{2\beta\varphi}dz^2$ with $\alpha^2=\frac1{12}$ for $n=1$.
\item P. G. O. Freund and M. A. Rubin, Phys. Lett. B 97 (1980) 233 — compactification on spheres by a form flux.
\item S. M. Carroll, J. Geddes, M. B. Hoffman and R. M. Wald, Phys. Rev. D 66 (2002) 024036 — stability of $dS_p\times S^q$ with flux and $\Lambda$.
\item L. Randall and R. Sundrum, Phys. Rev. Lett. 83 (1999) 3370 — the warped model RS1 ($2M^3R$, $V_{\rm hid}=-V_{\rm vis}=24M^3k$, $\Lambda=-24M^3k^2$).
\item W. D. Goldberger and M. B. Wise, Phys. Rev. Lett. 83 (1999) 4922 and Phys. Lett. B 475 (2000) 275 — stabilization by a bulk scalar; $f=\sqrt{24M^3/k}$, $m_\varphi^2=k^2v_v^2\epsilon^2e^{-2kr_c\pi}/(3M^3)$.
\item W. Israel, Nuovo Cim. B 44 (1966) 1 — junction conditions.
\item G. F. Giudice, R. Rattazzi and J. D. Wells, Nucl. Phys. B 595 (2001) 250 — Higgs–radion mixing ($\xi$, $Z$, $\theta$).
\item C. Csáki, M. L. Graesser and G. D. Kribs, Phys. Rev. D 63 (2001) 065002 — radion dynamics and phenomenology.
\item R. Camporesi and A. Higuchi, J. Geom. Phys. 20 (1996) 1 — the Dirac spectrum on spheres; Th. Friedrich, Math. Nachr. 97 (1980) 117 — the lower bound of Dirac eigenvalues.
\item J. G. Lee, E. G. Adelberger, T. S. Cook, S. M. Fleischer and B. R. Heckel, Phys. Rev. Lett. 124 (2020) 101101 — the Eöt-Wash bound $\lambda<38.6\,\mu$m for $|\alpha|=1$.
\item H. Weyl, Math. Ann. 71 (1912) 441 — the asymptotic distribution of eigenvalues; C. G. J. Jacobi (1829) — the four-square theorem.
\end{itemize}
\end{document}
